% Appendix C (continued): complete proofs for Exposés IV, V and VI, Definition IV.1 to Theorem VI.5.
% Sources: theory/propositions.json (statements, repaired proofs, public statements, round-2 notes),
% theory/round2_verdicts.json (suggested fixes for IV.2-IV.4, V.2-V.4, VI.3-VI.5, all applied), theory/framework.md
% (Exposés IV-VI), site/sections/07-lens.html, 08-base.html and 09-merge.html (refereed proofs). The restated results
% follow sections/04-lens.tex, 05-base.tex and 06-merge.tex. External results are Auxiliary results B.1-B.39.
% This file continues Appendix C and has no \section of its own.

% ====================================================================================================
\subsection{Definition IV.1}\label{proof:IV.1}

\begin{restate}{Definition}{IV.1}{forward and backward methods}
\begin{itemize}
  \item A \textbf{forward} method acts on the parameter port by a lens in the image of $\mathsf R$, namely
  $(\rho,\mathsf R[\rho])$.
  \item A \textbf{backward}, or cotangent-side, method acts by a lens whose forward part is $\id_\Theta$ and whose
  backward part is a \emph{compression} $c_t:T^*_\theta\Theta\to S$ into a state space $S$, paired with an
  \emph{anchor} $a_t:S\to T_\theta\Theta$. The update is
  \[\theta\leftarrow\theta-\eta\,a_t\big(\mathrm{opt}(c_t(\nabla L))\big).\]
\end{itemize}
Examples: GaLore, with $c_t=P_t\T$ and $a_t=P_t$ for $P_t$ from a gradient SVD; GaLore with a random projector, the same
with $P_t$ random; LISA, a layer-block coordinate projection paired with the same projection; MeZO, with
$c_t(g)=\langle z_t,g\rangle$ estimated by finite differences without a backward pass, and $a_t(\lambda)=\lambda z_t$.
Flora as published fits only as the trivial lens, with compression and anchor the identity and the random compression of
its momentum inside $\mathrm{opt}$.
\end{restate}

A definition needs no proof. This subsection fixes the conventions used in the proofs of Exposé~IV and checks the
identities that the definition and its table record.

\paragraph{Conventions}
As in \S\ref{proof:III.1}, weight space $\Theta=\R^{m\times n}$, or a product of such boxes, carries the Frobenius inner
product, covectors are written as matrices, and the transpose $\mathcal L^*$ of a linear map $\mathcal L$ is its adjoint.
We model a first-order optimizer as a machine with a state. At step $t$ it holds a state $\varsigma_t$, receives a tensor
$g_t$ (the gradient it is fed) and the tensor $x_t$ it updates, returns an output $U_t=F_t(\varsigma_t,g_t,x_t)$, moves to
the state $\varsigma_{t+1}=H_t(\varsigma_t,g_t,x_t)$, and the update is $x_{t+1}=x_t-\eta_tU_t$. The maps may depend on the
step counter $t$, for instance through Adam's bias corrections or Adafactor's relative step size. The optimizer
\emph{does not read the parameter tensor} when $F_t$ and $H_t$ do not depend on $x_t$. SGD, heavy-ball and Nesterov
momentum, Adam and AMSGrad, and Adafactor with \texttt{scale\_\allowbreak parameter=False} are of this kind; Adafactor's
factored second moments and its clipping of the update by $\mathrm{RMS}(U_t)$ read only $g_t$ and the state. The
following ingredients read $x_t$, or a gradient other than $g_t$:
\begin{itemize}
  \item coupled weight decay, which feeds $g_t+\lambda x_t$ in place of $g_t$;
  \item decoupled weight decay (AdamW with $\lambda>0$), which adds $-\eta_t\lambda x_t$ to the update;
  \item LAMB and LARS, which multiply the update by a trust ratio $\phi(\lVert x_t\rVert)/\lVert U_t\rVert$;
  \item Adafactor with parameter scaling, which multiplies the step by $\max(\epsilon_2,\mathrm{RMS}(x_t))$ (the option
  \texttt{scale\_\allowbreak parameter=\allowbreak True});
  \item clipping by the global norm of the full gradient, which rescales $g_t$ by $\min(1,\tau/\lVert G_t\rVert)$, where
  $G_t$ is the gradient of every trainable tensor of the procedure.
\end{itemize}

\begin{proof}[Verification of the identities in Definition IV.1]
\emph{The lens of a reparametrisation.} A smooth $\rho:Q\to\Theta$ has reverse derivative
$\mathsf R[\rho](q,\gamma)=J_\rho(q)\T\gamma=D\rho_q\T\gamma$ (\bgref{B.17}). For a composite $\rho\circ\sigma$ the chain rule
gives
\[\mathsf R[\rho\circ\sigma](q',\gamma)=D\sigma_{q'}\T D\rho_{\sigma(q')}\T\gamma=\mathsf R[\sigma]\big(q',\mathsf R[\rho](\sigma(q'),\gamma)\big),\]
which is composition of lenses, so $\rho\mapsto(\rho,\mathsf R[\rho])$ respects composition of reparametrisations. The
gradient that a forward method feeds its optimizer is the pullback $\nabla_q(L\circ\rho)=D\rho_q\T\nabla_\theta L=
\mathsf R[\rho](q,\nabla_\theta L)$.

\emph{A forward step to first order.} One step of a forward method sets $q\leftarrow q-\eta U$ with
$U=\mathrm{opt}(D\rho_q\T\nabla L)$, so the weight moves to
\[\rho(q-\eta U)=\theta-\eta\,D\rho_q\,U+O(\eta^2).\]
To first order in $\eta$ this is the backward update with $c=D\rho_q\T$ and $a=D\rho_q$. The two agree exactly when $\rho$ is
affine along the step, which is the situation of \thmref{Proposition}{IV.2}.

\emph{The examples.} For GaLore with left projection, $P_t\in\R^{m\times r}$ has orthonormal columns,
$S=\R^{r\times n}$, $c_t(G)=P_t\T G$ and $a_t(U)=P_tU$, so the update $W\leftarrow W-\eta\alpha P_t\,\mathrm{opt}(P_t\T G)$ is
the displayed one with $\alpha$ absorbed into $\eta$; with right projection, $c_t(G)=GP_t$ and $a_t(U)=UP_t\T$. In each case
$a_t=c_t^*$, since $\langle P_tU,G\rangle=\langle U,P_t\T G\rangle$. LISA's $S$ is the space of the sampled layers (with the
embedding and the head), $c_t$ restricts a gradient to it and $a_t$ includes it back, so again $a_t=c_t^*$. For MeZO,
$S=\R$, $a_t(\lambda)=\lambda z_t$ and $c_t(g)=\langle z_t,g\rangle=a_t^*(g)$. MeZO does not compute $c_t(\nabla L)$; it
estimates it by the central difference $\big(L(\theta+\varepsilon z_t)-L(\theta-\varepsilon z_t)\big)/(2\varepsilon)$,
whose error is $O(\varepsilon^2)$ for a loss with Lipschitz Hessian (\auxref{B.34}(a)), so it is the displayed update up to
its estimator. Flora as published (its reference implementation) returns $(1-\beta)\hat g_t+\beta\,\mathrm{up}_t(M_t)$, where $\hat g_t$ is the
Adafactor-normalised full gradient and $\mathrm{up}_t(M_t)$ the decompressed momentum \citep{hao2024flora}. The first term
is generically of full rank, so the update does not lie in $\operatorname{im}P_t$, and the only reading in the form of the
definition is $c_t=a_t=\id$ with the random compression inside $\mathrm{opt}$.

\emph{The cometric of a backward method.} With plain (S)GD on the compressed gradient, $\mathrm{opt}=\id$, the update is
$\theta\leftarrow\theta-\eta\,a_tc_t\nabla L$, so the induced cometric is $K_t=a_tc_t$ (\thmref{Definition}{III.1}). For the
first three examples $a_t=c_t^*$, so $K_t=a_ta_t^*$ is symmetric positive semidefinite; for GaLore with left projection,
$K_t(G)=P_tP_t\T G$.
\end{proof}

% ====================================================================================================
\subsection{Proposition IV.2}\label{proof:IV.2}

\begin{restate}{Proposition}{IV.2}{the degree-one collapse; known, \citet{xtorrobahennigen2025duality} for GaLore (their
Thm~1 and Cor.~3, App.~B for weight decay) and \citet{hao2024flora} for random projections}
Let $\rho_k(C)=W_k+\mathcal LC$ with $\mathcal L$ linear and $C_0=0$. Let $\mathrm{opt}$ be a first-order optimizer, with
no coupled or decoupled weight decay, whose state and output depend only on the sequence of gradients it is fed, and
which does not read the parameter tensor: SGD, heavy-ball or Nesterov momentum, Adam or AMSGrad, or Adafactor with
\texttt{scale\_\allowbreak parameter=False} (fixed or relative step, with its factorisation applied to the compressed
gradient in the same tensor shape as $C$). Run both procedures below with the same implementation of $\mathrm{opt}$, on
the same minibatches, and from the same optimizer state, fresh or carried over. Then they produce the same
$W$-trajectory over a period:
\begin{itemize}
  \item training $C$ with $\mathrm{opt}$ and merging at the end of the period;
  \item the backward method with $c=\mathcal L^*$ and $a=\mathcal L$, started at $W_k$.
\end{itemize}
AdamW with $\lambda>0$, LAMB and LARS trust ratios, parameter-scaled Adafactor (\texttt{scale\_\allowbreak
parameter=True}, the default in HF and in \texttt{galore\_torch}'s \texttt{GaLoreAdafactor}) and global-norm clipping of
the full gradient (such as HF Trainer's default \texttt{max\_grad\_norm=1.0}, whenever it triggers) all break the
identity. Injectivity of $\mathcal L$ is not needed.

In particular, on each matrix that GaLore projects (its other parameters train in full in both procedures), GaLore with
a projector $P_k$ fixed on each period equals one-sided LoRA $W_k+P_kC$ (frozen $B=P_k$, or frozen $A$ when projecting
on the right), merged and re-based every period, with $C$'s optimizer state (moments and step counter) carried over
unchanged across re-basings, with no reset. $P_k$ must be exactly GaLore's matrix, signs and column order included; its
span alone does not suffice, because the state is reused without transport. ReLoRA-style resets give a different
trajectory. The identity is exact with weight decay $0$, no full-gradient clipping, and an optimizer from the admissible
class; \texttt{galore\_torch}'s \texttt{GaLoreAdafactor} defaults to \texttt{scale\_\allowbreak parameter=True} and is
excluded. \texttt{GaLoreAdamW} is the transformers-style AdamW, which adds $\varepsilon$ to $\sqrt v$ before the bias
correction, so the LoRA side must use the same implementation. GaLore's scale $\alpha$ multiplies the optimizer output, so
it acts as a learning-rate factor; a LoRA-style scale inside $\rho$ gives a different trajectory under Adam with
$\varepsilon>0$.
\end{restate}

\begin{proof}
We use the optimizer model of \S\ref{proof:IV.1}. Fix a period of $T$ steps, write $L_t$ for the loss of the minibatch used
at step $t$ (the same in both procedures) and $\eta_t$ for the step size.

\emph{The two procedures.} In the first, the trainable tensor is $C$ with $C_0=0$ and the weight is
$W^{(1)}_t:=\rho_k(C_t)=W_k+\mathcal LC_t$. By the chain rule the gradient fed to $\mathrm{opt}$ is
\[g^{(1)}_t=\nabla_C(L_t\circ\rho_k)(C_t)=\mathcal L^*\nabla L_t\big(W^{(1)}_t\big),\]
$\mathrm{opt}$ returns $U^{(1)}_t$, and $C_{t+1}=C_t-\eta_tU^{(1)}_t$. In the second, $W^{(2)}_0=W_k$, the gradient fed to
$\mathrm{opt}$ is $g^{(2)}_t=c\big(\nabla L_t(W^{(2)}_t)\big)=\mathcal L^*\nabla L_t\big(W^{(2)}_t\big)$, $\mathrm{opt}$ returns
$U^{(2)}_t$, and $W^{(2)}_{t+1}=W^{(2)}_t-\eta_t\mathcal LU^{(2)}_t$.

\emph{Induction.} We show that for $0\le t\le T$ the weights $W^{(1)}_t$ and $W^{(2)}_t$ agree and the two optimizer
states agree. At $t=0$, $W^{(1)}_0=W_k+\mathcal L0=W_k=W^{(2)}_0$, and the states agree by hypothesis. Suppose both claims
hold at $t$. The two procedures evaluate the same minibatch loss at the same weight, so they compute the same weight
gradient and feed $\mathrm{opt}$ the same tensor $g_t:=g^{(1)}_t=g^{(2)}_t$. Since $\mathrm{opt}$ is the same
implementation and does not read the parameter tensor, its output and its next state are the same function of the state,
of $g_t$ and of $t$ in both procedures. Hence $U^{(1)}_t=U^{(2)}_t=:U_t$ and the next states agree. Then
\[W^{(1)}_{t+1}=W_k+\mathcal L(C_t-\eta_tU_t)=W^{(1)}_t-\eta_t\mathcal LU_t=W^{(2)}_{t+1}.\]
At the end of the period merging sets the weight to $\rho_k(C_T)=W^{(1)}_T=W^{(2)}_T$. The argument never inverts
$\mathcal L$, so injectivity is not needed. Every operation that $\mathrm{opt}$ applies to $g_t$ and its state, such as
Adafactor's factorisation in the shape of $C$ (which is the shape of $g_t$ in both procedures) and its update clipping, is
covered, since it reads nothing else.

\emph{The excluded ingredients.} Each of them reads a tensor that differs between the procedures, and in each case the
first step already separates the trajectories for suitable data; it suffices to exhibit this. Take $\mathcal L$ injective,
$W_k\ne0$ and a minibatch gradient $G_0$ with $\mathcal L^*G_0\ne0$.
\begin{itemize}
  \item \emph{Decoupled decay.} The first procedure sets $C_1=C_0-\eta_0U_0-\eta_0\lambda C_0=-\eta_0U_0$, so
  $W^{(1)}_1=W_k-\eta_0\mathcal LU_0$; the second sets $W^{(2)}_1=W_k-\eta_0\mathcal LU_0-\eta_0\lambda W_k$. They differ by
  $\eta_0\lambda W_k\ne0$. In general decay pulls $W$ towards $W_k$ in the first procedure and towards $0$ in the second.
  \item \emph{Coupled decay.} The first procedure feeds $\mathcal L^*G_0+\lambda C_0=\mathcal L^*G_0$, the second
  $\mathcal L^*(G_0+\lambda W_k)$. These differ whenever $\mathcal L^*W_k\ne0$, and for SGD the first steps then differ,
  since $\mathcal L$ is injective.
  \item \emph{Trust ratios and parameter scaling.} The first procedure reads $\lVert C_0\rVert=\mathrm{RMS}(C_0)=0$, the
  second $\lVert W_k\rVert>0$ and $\mathrm{RMS}(W_k)>0$. The resulting step sizes differ for generic $W_k$, because the
  first does not depend on $W_k$ and the second does.
  \item \emph{Global-norm clipping of the full gradient.} The first procedure clips by the norm of its trainable gradient
  $\mathcal L^*G_0$, the second (GaLore in HF Trainer) by the norm of the full weight gradient $G_0$. Under SGD with
  $\lVert\mathcal L^*G_0\rVert\le\tau<\lVert G_0\rVert$ the first step is $-\eta_0\mathcal L\mathcal L^*G_0$ in the first
  procedure and $-\eta_0(\tau/\lVert G_0\rVert)\mathcal L\mathcal L^*G_0$ in the second; they differ because
  $\mathcal L\mathcal L^*G_0\ne0$.
\end{itemize}

\emph{GaLore.} On a projected matrix with left projection, GaLore's update in period $k$ is
$W\leftarrow W-\eta\alpha P_k\,\mathrm{opt}(P_k\T G)$. This is the backward method with $c=P_k\T$ and $a=P_k$, and for
$\mathcal LC=P_kC$ we have $\mathcal L^*=P_k\T$; the factor $\alpha$ multiplies the output of $\mathrm{opt}$, so it is part
of the step size. With right projection, $\mathcal LC=CP_k\T$ and $\mathcal L^*G=GP_k$. Now induct over periods. Suppose that
at the start of period $k$ the whole model, every projected and unprojected tensor, is the same in GaLore and in the
re-based one-sided LoRA, and that so are all optimizer states. The unprojected parameters train in full with the same
optimizer in both, so they follow the same trajectory once the projected matrices do. GaLore computes $P_k$ from the
gradient at the start of the period, and the one-sided LoRA uses that same matrix as its frozen factor. The first part of
the proposition, applied with $\mathcal LC=P_kC$, gives the same trajectory through the period, and merging $P_kC_T$ into
$W_k$ gives the same $W_{k+1}$. GaLore keeps the state of $\mathrm{opt}$, which lives on tensors of the shape of $C$, and
the one-sided LoRA carries $C$'s state over while resetting $C$ to $0$; so the states at the start of period $k+1$ agree
too. This proves the corollary.

\emph{The exact matrix.} Replacing $P_k$ by $P_kO$ with $O\in\Ort(r)$ keeps the span and changes the trajectory when state
is carried. Take heavy-ball momentum with coefficient $\beta>0$ and SGD steps, $U_t=m_{t+1}=\beta m_t+g_t$, and a carried
momentum $m$ whose first row is nonzero, and let $O=\diag(-1,1,\dots,1)$. With $P_k$ the first step of the period moves
$W$ by $-\eta P_k(\beta m+P_k\T G)$. With $P_kO$ the compressed gradient is $OP_k\T G$, the momentum is reused without
transport, and $W$ moves by $-\eta P_kO(\beta m+OP_k\T G)=-\eta P_k(\beta Om+P_k\T G)$. The difference is
$-\eta\beta P_k(Om-m)\ne0$, since $P_k$ is injective and $Om-m$ is minus twice the first row of $m$. The same computation
applies to Adam's first moment.

\emph{Resets.} ReLoRA resets the optimizer state at every re-basing. For heavy-ball momentum the first step of period $k+1$
is then $-\eta P_{k+1}g$ instead of $-\eta P_{k+1}(\beta m+g)$, which differs whenever the carried momentum $m$ is
nonzero.

\emph{Implementation and scale.} With bias corrections $b_1=1-\beta_1^t$ and $b_2=1-\beta_2^t$, transformers-style AdamW
and \texttt{torch.optim.Adam} take the steps
\[\eta\,\frac{\sqrt{b_2}}{b_1}\,\frac{m_t}{\sqrt{v_t}+\varepsilon}\qquad\text{and}\qquad
\eta\,\frac{m_t/b_1}{\sqrt{v_t/b_2}+\varepsilon}=\eta\,\frac{\sqrt{b_2}}{b_1}\,\frac{m_t}{\sqrt{v_t}+\varepsilon\sqrt{b_2}} .\]
For $\varepsilon>0$ the two differ, because they use different effective $\varepsilon$; hence the hypothesis of one
implementation. Finally, put GaLore's $\alpha$ inside $\rho$, $\rho(C)=W_k+\alpha P_kC$. The compressed gradient becomes
$\alpha P_k\T G$, Adam's moments become $\alpha m_t$ and $\alpha^2v_t$, and the weight moves by
\[-\eta\,\alpha P_k\frac{\alpha\hat m_t}{\alpha\sqrt{\hat v_t}+\varepsilon}=-\eta\,\alpha P_k\frac{\hat m_t}{\sqrt{\hat v_t}+\varepsilon/\alpha},\]
against $-\eta\alpha P_k\,\hat m_t/(\sqrt{\hat v_t}+\varepsilon)$ for GaLore. These agree for $\varepsilon=0$, or when the LoRA
side's $\varepsilon$ is multiplied by $\alpha$, and differ for $\varepsilon>0$ and $\alpha\ne1$ whenever $\hat m_t\ne0$.
\end{proof}

\begin{remark*}{numerical checks}
The values reported after \thmref{Proposition}{IV.2} ($7\times10^{-16}$ over one period and $5\times10^{-15}$ over four
under Adam with the state carried; $8\times10^{-17}$ against the official \texttt{galore\_torch} optimizers; the
separations $0.029$, $0.020$, $0.030$ and $0.12$ for a reset, weight decay $0.1$, HF-default Adafactor and a flipped
column) instantiate the identity and the four breaking mechanisms above (Appendix~\ref{app:repro}).
\end{remark*}

% ====================================================================================================
\subsection{Proposition IV.3}\label{proof:IV.3}

\begin{restate}{Proposition}{IV.3}{integrability; after Frobenius' theorem}
Let $D$ be a smooth distribution of constant rank $k$ on an open set $U\subseteq\Theta$, for instance the anchor
distribution $D_\theta=\operatorname{im}a(\theta)$ of a time-independent backward method. Then $D$ is involutive on $U$ if
and only if through every $\theta\in U$ passes an injective immersion $\rho_\theta:(\R^k,0)\to(U,\theta)$ with
$\operatorname{im}D\rho_\theta=D$ along its image. A single such $\rho$ through $\theta_0$ gives involutivity only along
$\operatorname{im}\rho$. For example, $D=\spanop\{\partial_x,\ \partial_y+xz\,\partial_z\}$ on $\R^3$ has
$[\partial_x,\partial_y+xz\,\partial_z]=z\,\partial_z\notin D$ off $z=0$, yet the plane $z=0$ is an integral manifold.

The dynamic content needs separate hypotheses.
\begin{itemize}
  \item In continuous time, if $D$ is involutive, a backward flow with $\dot\theta\in D_\theta$ stays on the leaf through
  $\theta_0$.
  \item In the discrete update of \thmref{Definition}{IV.1}, the iterates stay on a leaf for all step sizes and all
  optimizer outputs $s$ (all losses) if and only if that leaf contains $(\theta+D_\theta)\cap U$ at every iterate, that
  is, if and only if the leaves are affine; an anchor constant over the period gives this (\thmref{Proposition}{IV.2}).
  \item Under gradient flow, dynamic equivalence with a forward lens $(\rho,\mathsf R[\rho])$ forces matching frames,
  $a_tc_t=D\rho\,\Lambda\,D\rho^\top$ along the trajectory. Hence $K=ac$ is symmetric positive semidefinite and, for
  injective $a$, $c=Ma^\top$ with $M$ symmetric positive semidefinite on $S$, possibly depending on $\theta$. If in
  addition $\rho$ is an immersion with $\dim Q=\rk D$, then $\rho$ is a local isometry from $(Q,\Lambda^{-1})$ onto the
  leaf with metric $(aMa^\top)^{-1}|_D$, so that metric must be flat. Forward methods with gauge fibres escape this
  requirement: $\rho(B)=BB^\top$ on symmetric matrices induces the curved Bures--Wasserstein metric, and balanced
  full-rank LoRA also induces a curved leaf metric. For discrete steps, exact agreement for every step size further
  forces $\rho$ to be affine along each horizontal line $q+t\,\Lambda D\rho_q^\top g$.
\end{itemize}

\emph{Instances.} Every projector-type backward method (GaLore with any period $T$ and any projector, its
random-projector variant, LISA, MeZO) has an anchor that does not depend on $\theta$ within its period. Its distribution
is constant, hence trivially involutive with affine leaves (vacuously at $T=1$, where a period is a single step). Under
the optimizer hypotheses of \thmref{Proposition}{IV.2} (no weight decay of either kind, no full-gradient clipping; GaLore
and MeZO need weight decay $0$, and LISA's AdamW with positive decay violates them), these methods are rebasing schemes
of linear forward methods with period $T$, including $T=1$, whose reach \thmref{Theorem}{V.3} describes. Each MeZO step is a step of the linear method
$\theta_t+\lambda z_t$, up to its SPSA estimator, and GaLore with $T=1$ is one-sided LoRA re-based every step. What
separates these methods from LoRA-type methods is the resampling schedule together with the carry-over of optimizer
state; integrability plays no part in it. Flora as published (Adafactor base, with a full-rank current-gradient term in
the code), APOLLO (a low-rank-estimated scaling applied to the full gradient) and Fira (a scaled full-rank residual) are
not of the form of \thmref{Definition}{IV.1} with their projections as anchors, because their updates leave
$\operatorname{im}a$.

Non-involutivity proper needs a $\theta$-dependent anchor. Full-batch GaLore with $T=1$, read as the distribution
$D_W=\operatorname{Hom}(\R^n,U_r(\nabla L(W)))$ (or $\operatorname{Hom}(V_r(\nabla L(W)),\R^m)$ when $m\ge n$, as GaLore
projects) wherever a singular-value gap exists, is involutive if and only if $U_r(\nabla L(W))$ (respectively
$V_r(\nabla L(W))$) is invariant along $D$. That fails generically for every $r\ge1$ once both matrix dimensions are at
least $2$. For every $r$ with $1\le r<\min(m,n)$ it fails for $L=\tfrac12\lVert W\rVert_F^2$, and numerically it fails
for least squares and cross-entropy with several samples. It holds, for instance, for single-sample least squares, and
automatically only when the distribution is a line field ($nr=1$). Minibatch GaLore with $T=1$ and MeZO violate
time-independence, because their anchors depend on step-wise randomness.
\end{restate}

\begin{proof}
Throughout, a \emph{leaf} of an involutive $D$ is a maximal connected integral manifold, and the leaves form a foliation
of $U$ (\auxref{B.31}(c)).

\emph{(1) The Frobenius equivalence.} Suppose $D$ is involutive and let $\theta\in U$. By \auxref{B.31}(b) there is a chart
$\varphi:V\to(-\delta,\delta)^{\dim\Theta}$ centred at $\theta$ in which $D$ is spanned by the first $k$
coordinate fields. Let $\psi:\R^k\to(-\delta,\delta)^k$ be the diffeomorphism $\psi_i(y)=\tfrac{2\delta}{\pi}\arctan y_i$ and
put $\rho_\theta(y)=\varphi^{-1}(\psi(y),0)$. Then $\rho_\theta(0)=\theta$, $\rho_\theta$ is an injective immersion into
$U$, and $\operatorname{im}D\rho_\theta(y)$ is spanned by the first $k$ coordinate fields at $\rho_\theta(y)$, which is
$D_{\rho_\theta(y)}$. Conversely, suppose such immersions exist through every point, and let $X,Y$ be sections of $D$ near
$\theta$. An immersion is locally an embedding (\auxref{B.1}(b)), so $\rho_\theta$ maps a neighbourhood of $0$ onto an
embedded integral manifold through $\theta$, and \auxref{B.31}(a) gives $[X,Y]_\theta\in D_\theta$. Since $\theta$ was
arbitrary, $D$ is involutive.

The same application of \auxref{B.31}(a) to the local pieces of $\operatorname{im}\rho$ shows that a single integral
immersion $\rho$ forces $[X,Y]\in D$ at the points of $\operatorname{im}\rho$; off $\operatorname{im}\rho$ it says nothing, as
the example shows. In the example,
$X_1=\partial_x$ and $X_2=\partial_y+xz\,\partial_z$ are linearly independent everywhere, so $D$ has constant rank $2$, and
$[X_1,X_2]=\partial_x(xz)\,\partial_z=z\,\partial_z$. A vector of $D$ at $(x,y,z)$ has the form $(a,b,bxz)$, which equals
$(0,0,z)$ only if $a=b=0$ and $z=0$; so $[X_1,X_2]\notin D$ wherever $z\ne0$, and $D$ is involutive on no open set meeting
$\{z\ne0\}$. On the plane $z=0$, $D=\spanop\{\partial_x,\partial_y\}$ is its tangent plane, so $\rho(u,v)=(u,v,0)$ is an
injective immersion with $\operatorname{im}D\rho=D$ along its image, and there $[X_1,X_2]=0\in D$.

\emph{(2) Continuous time.} Let $D$ be involutive, let $\theta(t)$, $t\in J$ an interval containing $0$, be an absolutely
continuous curve in $U$ with $\theta(0)=\theta_0$ and $\dot\theta(t)\in D_{\theta(t)}$ for almost every $t$, and let $N$ be
the leaf through $\theta_0$. Put $J_N=\{t\in J:\theta(t)\in N\}$, which contains $0$. Fix $t_*\in J$ and a flat chart around
$\theta(t_*)$ as in \auxref{B.31}(b). For $t$ near $t_*$ the curve stays in the chart and keeps its last $\dim\Theta-k$
coordinates constant, so it stays in the slice $\Sigma$ through $\theta(t_*)$. A slice is a connected integral manifold, and
a connected integral manifold that meets a leaf lies in it (\auxref{B.31}(c)), so if $\Sigma$ meets $N$ it lies in $N$. If $t_*\in J_N$, the slice $\Sigma$ meets $N$ at $\theta(t_*)$, so a neighbourhood of
$t_*$ lies in $J_N$: $J_N$ is open. If $t_n\to t_*$ with $t_n\in J_N$, then for large $n$ the point $\theta(t_n)$ lies in
$\Sigma\cap N$, so $\Sigma\subseteq N$ and $t_*\in J_N$: $J_N$ is closed. Hence $J_N=J$.

\emph{(3) Discrete steps.} The update is $\theta'=\theta-\eta\,a(\theta)s$ with $\eta>0$ and $s\in S$ arbitrary. As $\eta$
and $s$ range, $-\eta\,a(\theta)s$ ranges over the linear space $\operatorname{im}a(\theta)=D_\theta$, so the possible next
iterates in $U$ form $(\theta+D_\theta)\cap U$. They all lie on the leaf $N$ through $\theta$ iff $(\theta+D_\theta)\cap
U\subseteq N$; this is the first equivalence. Call the leaves \emph{affine} when every leaf $N$ equals $(\theta+D_\theta)\cap
U$ for each $\theta\in N$. If the containment holds at every point of every leaf, then for each $\theta\in N$ a small
$k$-disc of the plane $\theta+D_\theta$ around $\theta$ lies in $N$. Take a flat chart around $\theta$ (\auxref{B.31}(b)) and
shrink the disc into it. The leaves form a foliation (\auxref{B.31}(c)), so $N$ meets the chart in countably many slices, on
each of which the last $\dim\Theta-k$ coordinates are constant. On the connected disc these coordinates are continuous and
take countably many values, so they are constant, and the disc lies in the slice through $\theta$, a $k$-dimensional integral
manifold, in which it is open by invariance of domain. Hence its tangent planes, all equal to $D_\theta$, are the planes
$D_x$ at its points, so $D$ is
locally constant along $N$, so constant on the connected $N$; every path in $N$ then has velocity in $D_\theta$, so
$N\subseteq\theta+D_\theta$, and with the containment $N=(\theta+D_\theta)\cap U$. The converse is clear. For an anchor that is constant over the period, $D_\theta=V$ is defined
on $U=\Theta$, sections of $D$ are $V$-valued functions whose derivatives and brackets are again $V$-valued, and the leaves
are the planes $\theta+V$; the condition holds, as \thmref{Proposition}{IV.2} uses.

\emph{(4) Frames.} A backward method under gradient flow ($\mathrm{opt}=\id$) moves the weight by
$\dot\theta=-a_tc_t\nabla L$; the forward method moves it by $\dot\theta=-D\rho_q\Lambda D\rho_q\T\nabla L$
(\thmref{Definition}{III.1}). We read ``for every loss'' as including time-dependent losses, as minibatch training produces.
Let both methods follow the same trajectory up to time $t_1$ for a loss $L$, and switch at $t_1$ to $L+\langle\gamma,\cdot\rangle$.
The trajectories still agree up to $t_1$, and their right derivatives at $t_1$ agree, so
$a_{t_1}c_{t_1}(\nabla L+\gamma)=K_{q(t_1)}(\nabla L+\gamma)$ for every covector $\gamma$, that is,
$a_{t_1}c_{t_1}=D\rho\,\Lambda\,D\rho\T$ at $q(t_1)$. This operator is symmetric positive semidefinite, with range in
$\operatorname{im}a=D$. If $a$ is injective, let $a^+=(a\T a)^{-1}a\T$, so $P=aa^+$ is the orthogonal projector onto
$\operatorname{im}a$, and put $M=a^+(ac)(a^+)\T$, which is symmetric positive semidefinite. Since $ac$ is symmetric with
range in $\operatorname{im}a$, $ac=P(ac)=(ac)P$, so $ac=P(ac)P=aMa\T$, and $a(c-Ma\T)=0$ gives $c=Ma\T$.

Now let $\rho$ be an immersion with $\dim Q=\rk D=k$. Then $D\rho_q\Lambda D\rho_q\T$ has rank $k$ and range in $D$, so
$\operatorname{im}D\rho_q=D_{\rho(q)}$, and $\rho$ maps small open sets of $Q$ diffeomorphically onto open pieces of leaves.
On $D$ the cometric $K=aMa\T$ is invertible as a map from $D^*$ to $D$, and its inverse is a metric $g_D$ on $D$:
$g_D(v,v)=\langle\xi,v\rangle$ whenever $K\xi=v$. For $u\in T_qQ$ put $v=D\rho_qu$. Since $D\rho_q$ is injective,
$D\rho_q\T$ is onto, so there is $\xi$ with $D\rho_q\T\xi=\Lambda^{-1}u$, and then $K\xi=D\rho_q\Lambda\Lambda^{-1}u=v$. Hence
\[g_D(D\rho_qu,D\rho_qu)=\langle\xi,D\rho_qu\rangle=\langle D\rho_q\T\xi,u\rangle=\langle\Lambda^{-1}u,u\rangle .\]
So $\rho$ is a local isometry from $Q$ with the constant metric $\Lambda^{-1}$, which is flat, onto the leaf with metric
$g_D=(aMa\T)^{-1}|_D$, and by \auxref{B.32}(a) the leaf metric is flat.

\emph{(5) Gauge fibres: $BB\T$.} Let $\rho(B)=BB\T$ on $Q=\R^{n\times n}$ with $\Lambda=I$, landing in $\Sym(n)$, and pair
covectors with $\Sym(n)$ by the trace. Then $D\rho_B(\dot B)=\dot BB\T+B\dot B\T$, and for symmetric $G$,
$\langle G,\dot BB\T+B\dot B\T\rangle=\langle2GB,\dot B\rangle$, so $D\rho_B\T G=2GB$ and
\[K(G)=D\rho_B(2GB)=2(G\theta+\theta G),\qquad\theta=BB\T .\]
This depends only on $\theta$, not on the point of the fibre $\{BR:R\in\Ort(n)\}$. On the positive definite matrices it is
invertible (\auxref{B.20}), so $D=\Sym(n)$ there. Put $a=K^{1/2}$ and $c=a\T$. The backward flow
$\dot\theta=-K(\theta)\nabla L$ and the forward flow pushed to $\theta$ solve the same ordinary differential equation with
smooth right-hand side, so they agree for every loss. The leaf metric $K^{-1}$ is a constant multiple of the
Bures--Wasserstein metric (\auxref{B.32}(b)), which is not flat for $n\ge2$. Here $\dim Q=n^2>\dim\Sym(n)$, so $\rho$ is no
immersion, in line with (4).

\emph{(6) Gauge fibres: balanced full-rank LoRA.} Take $m=n=r\ge2$, $s=1$ and equal block rates $\eta_A=\eta_B=1$, and
restrict to balanced factors, $B\T B=AA\T$, with $W=BA$ invertible; gradient flow preserves balancedness by the Noether
charge of \thmref{Theorem}{III.5}. From $B\T B=AA\T$, $W\T W=A\T(B\T B)A=(A\T A)^2$ and $WW\T=B(AA\T)B\T=(BB\T)^2$, so
$A\T A=\lvert W\rvert:=(W\T W)^{1/2}$ and $BB\T=\lvert W\T\rvert:=(WW\T)^{1/2}$, the unique positive definite square roots.
By \thmref{Observation}{III.2},
\[K_W(G)=G\,\lvert W\rvert+\lvert W\T\rvert\,G ,\]
a function of $W$ alone, invertible by \auxref{B.20}; as in (5) a backward method with $a=K^{1/2}$, $c=a\T$ reproduces the
flow, and the leaf is an open subset of $\GL_n$ with metric $g_W(X,X)=\langle K_W^{-1}X,X\rangle$. The transpose $\tau(W)=W\T$
is an isometry of $g$: $K_{W\T}(G\T)=(K_WG)\T$, so $K_{W\T}^{-1}(X\T)=(K_W^{-1}X)\T$ and
$g_{W\T}(X\T,X\T)=g_W(X,X)$. Its fixed points in $\GL_n$ are the invertible symmetric matrices, and the positive definite ones
form an open subset $F$ of them, an embedded submanifold on which $d\tau=\id$. A $g$-geodesic $\gamma$ with $\gamma(0)\in F$ and $\dot\gamma(0)$ tangent to $F$
is carried by $\tau$ to a geodesic with the same initial data (\auxref{B.38}(a)), so $\tau\circ\gamma=\gamma$ and
$\gamma$ stays in $F$ for small times. Every geodesic tangent to $F$ stays in $F$, so $F$ is totally geodesic, and by the
Gauss equation its intrinsic sectional curvatures equal those of $g$ on its tangent planes (\auxref{B.38}(b)). On $F$,
$\lvert W\rvert=\lvert W\T\rvert=W$, and the induced metric is $\langle L_W^{-1}X,X\rangle$ with $L_W(G)=GW+WG$ on $\Sym(n)$,
half the cometric of (5); this is a constant multiple of the Bures--Wasserstein metric, which has a plane of positive
sectional curvature for $n\ge2$ (\auxref{B.32}(b)). So $g$ has a plane of nonzero sectional curvature and is not flat.

\emph{(7) Discrete steps and horizontal lines.} A plain GD step of the forward method from $q$, with $g=\nabla L(\rho(q))$,
lands at $\rho(q-\eta v)$ with $v=\Lambda D\rho_q\T g$; the backward step lands at $\rho(q)-\eta\,a c\,g$. Agreement for small
$\eta$ forces $ac=D\rho_q\Lambda D\rho_q\T$ as in (4), so the backward step is $\rho(q)-\eta\,D\rho_qv$. Agreement for every
$\eta>0$, and for $-g$ as well as $g$ (another loss), says $\rho(q+tv)=\rho(q)+t\,D\rho_qv$ for all real $t$ in the domain:
$\rho$ is affine along the line $q+t\Lambda D\rho_q\T g$.

\emph{(8) Instances.} Within a period, GaLore's anchor is a fixed matrix $P_k$, LISA's a fixed coordinate inclusion, and
MeZO's the fixed vector $z_t$ for its single step; none depends on $\theta$, so (3) applies on $U=\Theta$ (vacuously when a
period is one step). Each pair has $c=a^*$ (\S\ref{proof:IV.1}), so with $\mathcal L=a$ it is the backward method of
\thmref{Proposition}{IV.2}, and under that proposition's optimizer hypotheses each period is a period of the linear forward
method $\theta_k+\mathcal LC$, merged at its end. For MeZO, $\mathcal L(\lambda)=\lambda z_t$, and the identification holds up
to the error of the central difference (\auxref{B.34}). Integrability holds for all of them, so it cannot tell them apart,
and by \thmref{Proposition}{IV.2} what distinguishes them from LoRA is how the projector changes and which state is
carried. Flora's update contains the full-rank term $(1-\beta)\hat g_t$, APOLLO's is $G\,\diag(s)$ with $G$ the full
gradient \citep{zhu2024apollo}, and Fira's adds a scaled multiple of $(I-P_tP_t\T)G$ \citep{chen2024fira}; for gradients with
$(I-P_tP_t\T)G\ne0$ each update leaves $\operatorname{im}P_t$, so none is of the displayed form with $a=P_t$.

\emph{(9) The GaLore field.} Let $L$ be smooth. On the open set where $\sigma_r(\nabla L(W))>\sigma_{r+1}(\nabla L(W))$,
let $\Pi(W)$ be the orthogonal projector onto $U_r(\nabla L(W))$, the top-$r$ eigenspace of $S(W)=\nabla L(W)\nabla L(W)\T$. It
is a smooth function of $W$ there, being a Riesz projector of $S(W)$ whose spectral gap persists (\auxref{B.37}). The fields
$X_E(W)=\Pi(W)E$, $E\in\R^{m\times n}$, span $D_W=\{Z:\operatorname{col}Z\subseteq U_r\}$. Write $\Pi'(Z)=D\Pi(W)[Z]$.
Differentiating $\Pi^2=\Pi$ gives $\Pi'\Pi+\Pi\Pi'=\Pi'$, hence $\Pi\Pi'\Pi=0$ and $(I-\Pi)\Pi'(I-\Pi)=0$. With
$N(Z):=(I-\Pi)\Pi'(Z)\Pi$ and the symmetry of $\Pi'(Z)$,
\begin{equation}\label{eq:C3-proj-deriv}
\Pi'(Z)=N(Z)+N(Z)\T ,
\end{equation}
so $\Pi'(Z)=0$ iff $N(Z)=0$. The bracket is $[X_E,X_F]=\Pi'(\Pi E)F-\Pi'(\Pi F)E$, and its component normal to $D$ is
\[(I-\Pi)[X_E,X_F]=N(\Pi E)\,\Pi F-N(\Pi F)\,\Pi E .\]
\emph{Criterion.} If $U_r$ is invariant along $D$, that is $\Pi'(Z)=0$ for $Z\in D$, the normal component vanishes and $D$ is
involutive. Conversely, let $D$ be involutive and $n\ge2$. For $u,u'\in U_r$ and $x,y\in\R^n$ linearly independent, take
$\Pi E=ux\T$ and $\Pi F=u'y\T$: involutivity gives $N(ux\T)u'\,y\T=N(u'y\T)u\,x\T$. The left side has its rows in
$\spanop\{y\}$ and the right side in $\spanop\{x\}$, so both vanish, and $N(ux\T)u'=0$. As $u'$ ranges over $U_r$ and
$N(Z)$ vanishes on $U_r^\perp$, $N(ux\T)=0$ for all $u\in U_r$ and $x\in\R^n$ (given $x$, pick $y$ independent of it), so
$N=0$ on $D$ by linearity, and $U_r$ is invariant along $D$. For right projection apply this to the transpose. The failure proved next covers every $1\le r<\min(m,n)$
with $m,n\ge2$; the remaining case with a gap is $n=r=1$, where $D_W=U_1$ is a line field, and a line field is always
involutive, since $[fX,gX]=(f\,Xg-g\,Xf)X$.

\emph{Failure for $L=\tfrac12\lVert W\rVert_F^2$.} Here $\nabla L(W)=W$. Let $p=\min(m,n)\ge2$, $1\le r<p$, and
$W=\sum_{i\le p}\sigma_ie_ie_i\T$ with $\sigma_1>\dots>\sigma_p>0$, so $U_r=\spanop\{e_1,\dots,e_r\}$ and $S=WW\T$ is
diagonal with entries $\lambda_i=\sigma_i^2$ ($i\le p$) and $0$. For a smooth symmetric family $S(t)$ with $S(0)=S$,
differentiating $[S(t),\Pi(t)]=0$ gives $[\dot S,\Pi]+[S,\dot\Pi]=0$; the $(i,j)$ entry with $i\le r<j$ reads
$-\dot S_{ij}+(\lambda_i-\lambda_j)\dot\Pi_{ij}=0$, and $\dot\Pi$ vanishes on the diagonal blocks by the identities above.
So $\dot\Pi_{ij}=\dot S_{ij}/(\lambda_i-\lambda_j)$ for $i\le r<j$. Take $Z=e_1x\T\in D$, so $\dot S=ZW\T+WZ\T$. Then
$\dot S_{1j}=x\T W\T e_j=\sigma_jx_j$ for $r<j\le p$, $\dot S_{1j}=0$ for $j>p$, and $\dot S_{ij}=0$ for $2\le i\le r<j$. Hence
\[N(e_1x\T)=\sum_{j=r+1}^{p}\frac{\sigma_jx_j}{\sigma_1^2-\sigma_j^2}\,e_je_1\T .\]
With $E=e_1e_{r+1}\T$ and $F=e_1e_1\T$ we get $N(\Pi F)=0$ and
\[(I-\Pi)[X_E,X_F]=N(e_1e_{r+1}\T)\,e_1e_1\T=\frac{\sigma_{r+1}}{\sigma_1^2-\sigma_{r+1}^2}\,e_{r+1}e_1\T\ne0 ,\]
so $D$ is not involutive at $W$. At $m=n=2$, $r=1$ and $W=\diag(2,1)$, the frame $X_i=X_{e_1e_i\T}$ gives the full bracket
$[X_1,X_2]=\Pi'(e_1e_1\T)e_1e_2\T-\Pi'(e_1e_2\T)e_1e_1\T=0-\tfrac13(e_1e_2\T+e_2e_1\T)e_1e_1\T=-\tfrac13e_2e_1\T$; the frame
$X_i=u(W)e_i\T$, with $u$ the top left singular vector, gives the same value, since $Du[e_1e_1\T]=0$ and
$Du[e_1e_2\T]=\tfrac13e_2$ at this $W$.

\emph{Genericity.} For least squares $L(W)=\tfrac12\lVert WX-Y\rVert_F^2$ the normal component of the bracket of two fixed
frame fields is a real-analytic function of $(W,X,Y)$ on the open set where the gap exists. It does not vanish identically
on the connected component containing $(W,I_n,0)$ with $W$ as above, where it is the value just computed, so it is nonzero
on an open dense subset of that component (\auxref{B.10}). The values for cross-entropy and for several samples are the
numerical checks reported with the proposition.

\emph{Single-sample least squares.} With one sample, $\nabla L(W)=(Wx-y)x\T$ has rank one and, where $Wx\ne y$,
$U_1=\spanop\{Wx-y\}$. Moving $W$ along $Z=uv\T$ with $u\in U_1$ changes $Wx-y$ by $u(v\T x)\in U_1$, so $U_1$ is constant
along $D$, and $D$ is involutive by the criterion.

\emph{Time-dependence.} Minibatch GaLore with $T=1$ draws a new minibatch, and MeZO a new $z_t$, at every step, so their
anchors are functions of $(t,\theta)$ and not of $\theta$ alone.
\end{proof}

% ====================================================================================================
\subsection{Observation IV.4}\label{proof:IV.4}

\begin{restate}{Observation}{IV.4}{backward reach and the backprop cone; an upper bound}
Assume reverse mode in which each box keeps its own residual, with no checkpointing across boxes, and take boxes at the
granularity where adapters are separate boxes (for LoRA-FA, $A_0$ and $B$ separately). A \emph{VJP residual} is an input,
an output, a sufficient statistic such as a ReLU mask, or a recipe to recompute one of these inside the box. Then it
suffices to keep a VJP residual for
\begin{itemize}
  \item each trainable box whose parameter-VJP depends on its input: linear and multiplicative boxes such as $W$, $A$, $B$
  and (IA)$^3$ scales; additive boxes such as biases and prompts need nothing;
  \item each \emph{nonlinear} box, bilinear boxes such as the attention products included, on a directed path from a
  trainable box to the loss.
\end{itemize}
Deterministic frozen linear boxes, and all boxes off this cone, need nothing; stochastic linear boxes on the cone
(dropout) need their mask or random-number state. Alternatively, the $k$ coordinates of $P^\top\nabla L$ can be computed
exactly by $k$ forward-mode JVPs, with no reverse pass and $O(1)$-layer activation memory, each JVP costing a small
constant multiple (about $1.5$ to $3$) of a forward pass. This always uses less activation memory than un-checkpointed
reverse mode, but it is time-competitive with one VJP (about $3$ forward passes) only for $k$ of about $1$ or $2$. It also
needs $P$ before the passes, so it cannot produce GaLore's SVD refresh, where $P$ depends on $\nabla L$, more cheaply than a
backward pass ($\nabla L$ itself would take $\lvert\theta\rvert$ JVPs).
\end{restate}

\begin{proof}
\emph{The computation.} The network is a directed acyclic graph of boxes. A box receives inputs $x$ (activations), may carry
parameters $p$, and produces an output $y=f(x,p)$; the loss is the output of the last box. Reverse mode propagates
cotangents backwards. For a box with output cotangent $\bar y$ it computes the input cotangents $\bar x=J_xf(x,p)\T\bar y$
and, for a trainable box, the parameter cotangent $\bar p=J_pf(x,p)\T\bar y$ (\bgref{B.17}). Call the \emph{cone} the set of
boxes on a directed path from a trainable box to the loss.

\emph{Pruning.} The parameter gradients are sums over paths from the parameters to the loss of products of Jacobians, so
only cotangents of wires on such paths contribute. A wire whose value does not depend on any trainable parameter carries a
cotangent that is never multiplied into a parameter gradient, so its cotangent need not be formed. Hence it suffices to
compute $\bar y$ for the outputs of boxes on the cone, $\bar x$ for those inputs of cone boxes that are produced on the
cone, and $\bar p$ for the trainable boxes. Boxes off the cone take no part in the backward sweep and need no residual.

\emph{What each box on the cone needs.} We list, for each kind of box, what the backward sweep reads besides $\bar y$ and
the weights, which are stored anyway.
\begin{itemize}
  \item A deterministic frozen linear or affine box, $y=Wx+b$: $\bar x=W\T\bar y$ does not involve $x$, so no residual.
  \item A dropout box, $y=\mathsf m\odot x/(1-p)$: $\bar x=\mathsf m\odot\bar y/(1-p)$, so it needs its mask, or the random
  state that regenerates it.
  \item A trainable linear box, $y=Wx$: $\bar W=\bar y\,x\T$, summed over tokens, so it needs its input $x$; $\bar x=W\T\bar
  y$ needs nothing more. A trainable multiplicative box, $y=\ell\odot x$: $\bar\ell=\sum\bar y\odot x$, so it needs $x$,
  and $\bar x=\ell\odot\bar y$.
  \item A trainable additive box, $y=x+b$, or a prompt, $y=[P;x]$: $\bar b=\sum\bar y$ and $\bar P$ is the slice of $\bar y$
  at the prompt positions, and $\bar x$ is $\bar y$ or its slice; no residual.
  \item A nonlinear box, $y=f(x)$, possibly with frozen parameters: $\bar x=J_f(x)\T\bar y$ depends on $x$ through
  $J_f(x)$, so it needs something from which $J_f(x)\T\bar y$ can be formed: the input, the output (for tanh, sigmoid or
  softmax), a mask (for ReLU), or a recipe to recompute one of them. A bilinear box is nonlinear in its joint input; for
  $y=a\odot b$, $\bar a=\bar y\odot b$ and $\bar b=\bar y\odot a$, and for the attention products $QK\T$ and $PV$ the same
  holds with matrix products.
\end{itemize}
A trainable nonlinear box lies on the cone, being on a path from itself to the loss, so it is covered by the second bullet,
and also by the first if its parameter-VJP reads its input. Processing the cone in reverse topological order, each box has
$\bar y$ available (it is a sum of input cotangents of its readers on the cone) together with its residual, so the sweep
computes every parameter gradient exactly. This proves that the residuals listed suffice. For LoRA-FA, $A_0$ is a frozen
linear box and needs nothing, while $B$ is a trainable linear box whose input is $A_0x$, so $B$ keeps $A_0x$.

\emph{The forward-mode route.} The $i$-th coordinate of $P\T\nabla L$ is $\langle p_i,\nabla L(\theta)\rangle=DL(\theta)[p_i]$,
the directional derivative along the $i$-th column of $P$. By \auxref{B.33}(a), forward mode computes it alongside one
evaluation of $L$ at a small constant multiple of that cost, carrying each tangent with its value and discarding both once
consumed. In a layered network this keeps the values and tangents of $O(1)$ layers alive at a time, whereas un-checkpointed
reverse mode keeps the residuals of every box on the cone until the backward sweep reaches it. The $k$ coordinates cost
$k$ such passes. One reverse pass costs about three forward passes, so with a per-JVP cost of $1.5$ to $3$ forward passes the
forward route is no slower only when $1.5k$ to $3k$ is at most about $3$, that is, for $k$ of about $1$ or $2$. The
directions $p_i$ must be fixed before the passes. GaLore's refresh computes $P$ from the SVD of the full gradient $\nabla L$
of each projected matrix, and $\nabla L$ by forward mode costs one pass per coordinate, $\lvert\theta\rvert$ passes, against
one reverse pass (\auxref{B.33}(b)).

\emph{MeZO.} The central difference estimates $DL(\theta)[z]$ with error $O(\varepsilon^2)$ when $L$ has a Lipschitz
Hessian, and is unbiased for the gradient of the Gaussian smoothing of $L$ (\auxref{B.34}).
\end{proof}

% ====================================================================================================
\subsection{Definition V.1}\label{proof:V.1}

\begin{restate}{Definition}{V.1}{base change and rebasing}
\begin{itemize}
  \item[(a)] A \emph{function-preserving} pointed reparametrisation $\phi:(\Theta',f',\theta_0')\to(\Theta,f,\theta_0)$,
  with $f\circ(\phi\times\id_X)=f'$, induces by postcomposition
  \[\phi_!:\Meth(\Theta',\theta_0')\to\Meth(\Theta,\theta_0),\qquad(Q,q_0,\rho)\mapsto(Q,q_0,\phi\circ\rho).\]
  \item[(b)] A map $\kappa:\Theta\to\Theta$ that does \emph{not} preserve the function, such as a quantiser, moves the
  base point. Its \textbf{pointing defect} is $e=\kappa\theta_0-\theta_0$.
  \item[(c)] A \textbf{rebasing scheme} (merge-and-restart) picks $\theta_{k+1}\in\Im M(\theta_k)$. Its reachable sets
  are $R_K$, the points $\theta_K$ it can reach in $K$ cycles, and $R_\infty=\bigcup_KR_K$.
\end{itemize}
\end{restate}

A definition needs no proof. We check that $\phi_!$ is a functor that preserves the realised functions, and record the
recursion for the reachable sets that the proofs of \thmref{Theorem}{V.3} use.

\begin{proof}[Verification of the claims in Definition V.1]
(a) Since $\phi(\theta_0')=\theta_0$, the map $\phi\circ\rho$ sends $q_0$ to $\theta_0$, so $\phi_!M$ is an object of
$\Meth(\Theta,\theta_0)$. A morphism $h:M\to N$ is a pointed smooth map with $\rho_N\circ h=\rho_M$
(\thmref{Definition}{I.2}); then $(\phi\circ\rho_N)\circ h=\phi\circ\rho_M$, so the same $h$ is a morphism
$\phi_!M\to\phi_!N$. Setting $\phi_!h=h$ preserves identities and composition, so $\phi_!$ is a functor, the
postcomposition functor of slice categories (\bgref{A.19}). It changes neither the parameter space nor the realised
functions: for $q\in Q$,
\[f\big(\phi(\rho(q)),x\big)=f'\big(\rho(q),x\big)\qquad\text{for all }x\in X,\]
so $\phi_!M$ computes exactly the functions $M$ computes, and $\Im\phi_!M=\phi(\Im M)$.

(b) is a definition. (c) Put $R_0=\{\theta_0\}$. Each cycle trains the method $M(\theta_k)$ over the current base and
merges, so the points reachable in $K+1$ cycles are
\[R_{K+1}=\bigcup_{\theta\in R_K}\Im M(\theta).\]
When every $M(\theta)$ is pointed at $\theta$, that is $\rho_\theta(q_0)=\theta$, we have $\theta\in\Im M(\theta)$, so
$R_K\subseteq R_{K+1}$ and $R_\infty$ is an increasing union.
\end{proof}

% ====================================================================================================
\subsection{Proposition V.2}\label{proof:V.2}

\begin{restate}{Proposition}{V.2}{quantisation and splitting}
Write $W$ for the full-precision weight, $\kappa$ for the quantiser, $e=\kappa W-W$, $\delta_0=\delta(q_0)$, and $s$ for the
LoRA scale $\alpha/r$ ($\alpha/\sqrt r$ under rsLoRA).

(a) Additive families commute with translation, $M(\theta+v)=\tau_v\circ M(\theta)$. Hence QLoRA is LoRA pointed at
$\kappa\theta_0$, with defect $e$, and it needs no new weight-space structure. Since $-e$ is generically of full rank,
$\theta_0\notin\Im\mathrm{LoRA}_r(\kappa\theta_0)$ unless $\rk e\le r$.

(b) For splitting families derived from an SVD or from data, the order of operations matters, and three orders occur in
practice.
\begin{itemize}
  \item \emph{Split, then quantise the residual} (QPiSSA; HF's recommended PiSSA and CorDA workflow; the quantisation
  half-step of LoftQ for $t\ge2$, with $W_{\rm res}=W-\delta(q_{t-1})$). The defect is $\kappa(W_{\rm res})-W_{\rm res}$.
  \item \emph{Quantise, then fit the error} (LoftQ with one iteration, HF's default \texttt{loftq\_iter=1}; HF
  \texttt{replace\_\allowbreak lora\_\allowbreak weights\_\allowbreak loftq}). HF sets the factors from
  $\operatorname{SVD}_r(W-\kappa W)$ without dividing by $s$, unlike its PiSSA and OLoRA initialisations, so
  $\delta_0=-s\operatorname{SVD}_r(e)$ and the defect is $e-s\operatorname{SVD}_r(e)$, with
  \[\lVert e-s\operatorname{SVD}_r(e)\rVert^2=\lVert e-\operatorname{SVD}_r(e)\rVert^2+(1-s)^2\lVert\operatorname{SVD}_r(e)\rVert^2.\]
  The defect is $e$ minus its best rank-$r$ approximation (of rank $\rk e-r$) iff $s=1$, which is HF's bare default
  $r=\alpha=8$. For $s=2$ its norm equals $\lVert e\rVert$, no better than QLoRA. With a callback,
  \texttt{replace\_\allowbreak lora\_\allowbreak weights\_\allowbreak loftq} can roll a module back to the zero
  initialisation, whose defect is $e$.
  \item \emph{Quantise, split, re-quantise} (HF \texttt{olora\_init} on bitsandbytes weights). The defect is
  $e+\kappa(r')-r'$ with $r'=\kappa W-\delta_0$.
\end{itemize}
Quantise-then-split has defect exactly $e$ iff the frozen part equals $\kappa W-\delta_0$, for instance $\kappa W$ kept
quantised minus a frozen copy of the rank-$r$ $\delta_0$, at a cost of about $r(m+n)$ extra numbers. Re-quantising the
residual, as HF \texttt{olora\_init} does, generically adds $\kappa(r')-r'$.

Action families whose pointing does not depend on the base (OFT at $R=I$, BOFT at $R=I$ with unit scale, HiRA at $B=0$),
evaluated at $\kappa\theta_0$, simply carry the defect $e$. Methods whose pointing is computed from the base have an order
ambiguity of their own. DoRA's magnitude initialised from $\kappa\theta_0$ gives defect $e$, while initialised from
$\theta_0$ it gives $\diag\big(\lVert w_{0,i}\rVert/\lVert(\kappa\theta_0)_i\rVert\big)\,\kappa\theta_0-\theta_0$. For
action families a base change also moves the absolute invariants, along with the base point. NF4's exact zero code
enlarges the zero set that HiRA must preserve, and exactly orthogonal (exact-Cayley) OFT preserves the Gram matrix of
$\kappa\theta_0$ rather than that of $\theta_0$; HF's default Cayley--Neumann OFT preserves it only approximately.

(c) LoftQ alternates two updates on $\lVert\theta_0-Q-LR\rVert$, an exact rank-$r$ SVD step and a round-to-nearest
quantisation step with blockwise scales recomputed from its input. The second is not an exact projection, so the
objective need not decrease monotonically. With $s=1$ the output is pointed up to the attained residual. It is then an
\emph{approximate section}, a choice $(Q,q_0)$ with reduced $\lVert\rho(q_0)-\theta_0\rVert$. In general, with $Q_T$ and
the low-rank fit $L_TR_T$ after $T$ iterations, HF's deployed defect is $(Q_T+L_TR_T-W)+(s-1)L_TR_T$, because HF's internal
residual also omits $s$. The official LoftQ checkpoints use $r=64$ and $\alpha=16$, so $s=1/4$, and then the deployed
defect can grow with $T$ while the attained residual falls.

(d) For a generic high-precision adapter $\Delta$, fine-tune-then-quantise and quantise-then-fine-tune do not commute. The
first lands on the quantisation grid. In the second, the exactly merged weight $\kappa\theta_0+\Delta$ leaves the
quantised type. For a uniform fixed-scale INT-$k$ grid, $\kappa\theta_0+\Delta$ stays on the grid iff $\Delta$ lies on
the scale lattice (within range), so there the obstruction is the off-grid update. NF4 is non-uniform and not closed under
addition, and absmax grids whose scales are recomputed are not closed either, so for QLoRA's NF4 the grid itself is part of
the obstruction. Type-preserving merges, such as HF PEFT's dequantise--add--requantise for bitsandbytes 4- and 8-bit
layers, re-quantise with recomputed block scales, so they stay in the quantised type but are lossy. Under affine group
quantisation QA-LoRA is the structured exception. Its group-pooled adapter is constant on each quantisation group and
merges into the zero-points, so the merged model stays INT-$k$.
\end{restate}

\begin{proof}
\emph{Quantisers.} We use two properties of the blockwise quantisers in practice (NF4 and INT-$k$, with a fixed scale or
the absmax scale of each block, \citealp{dettmers2023qlora,dettmers20218}). First, $\kappa$ rounds each entry, divided by
its block's scale, to the nearest point of a finite codebook and multiplies back, so its values, the \emph{representable}
matrices, form a Lebesgue-null set: a finite set for fixed scales, and for absmax scales a finite union of sets with one
free scale per block of $\beta\ge2$ entries. Second, $\kappa$ is idempotent, $\kappa(\kappa X)=\kappa X$, and it fixes
exactly the representable matrices.

(a) An additive family has $\rho_\theta(q)=\theta+\delta(q)$ with $\delta$ independent of the base and $\delta(q_0)=0$
(\thmref{Definition}{I.2}). Then $\rho_{\theta+v}(q)=v+\rho_\theta(q)=\tau_v(\rho_\theta(q))$, which is the first claim.
QLoRA trains zero-initialised LoRA over the frozen base $\kappa\theta_0$, so it is $\mathrm{LoRA}_r(\kappa\theta_0)$, with
$\rho(q_0)=\kappa\theta_0$ and defect $\rho(q_0)-\theta_0=e$. Since every matrix of rank at most $r$ has the form $sBA$
($s\ne0$), $\theta_0=\kappa\theta_0+sBA$ for some $(B,A)$ iff $\rk(-e)\le r$.

\emph{The word ``generically''.} For a fixed-scale quantiser, $\kappa$ is constant on each open cell of a decomposition of
$\Theta$ into countably many cells and a null set of boundaries, so on a cell $e=K-W$ with $K$ constant. The rank-deficient
matrices are the common zeros of the maximal minors, a proper algebraic subset, and $W\mapsto K-W$ is a bijection, so $e$
has full rank for almost every $W$ in each cell, hence for almost every $W$ (\auxref{B.10}). For an absmax quantiser with
blocks of $\beta\ge2$ consecutive entries of a row, take the cells on which every code and the position of every block
maximum are constant. With the scale stored exactly, the block maximum is represented exactly, so $e$ vanishes at those positions, and each other
entry is $e_{ij}=a\,c_{ij}-W_{ij}$ with $a$ the scale of its block and $c_{ij}$ its code; varying $W_{ij}$ alone moves
$e_{ij}$ alone. So on each cell $e$ fills an open subset of the space $Z$ of matrices vanishing at the block maxima, and its
rank is, for almost every $W$ in the cell, the largest rank in $Z$. A row of length $n$ has at most $\lceil n/\beta\rceil+1$
forced zeros, so $Z$ contains a matrix of rank $r+1$ whenever $r+1\le n-\lceil n/\beta\rceil-1$ and $r<m$ (place $r+1$
nonzero entries in distinct rows and columns, avoiding the forced zeros, greedily). Hence $\rk e>r$, and
$\theta_0\notin\Im\mathrm{LoRA}_r(\kappa\theta_0)$, for almost every $W$ in this range. Full rank holds for almost every
$W$ in each cell whose zero pattern admits a matrix of full rank; the other cells, in which the block maxima of many rows
fall in a few common columns, have positive but tiny measure (for instance, a column of zeros needs the maximum of its
block in every row, probability $\beta^{-m}$ per column for weights with independent continuous entries).

(b) In each pipeline the initial weight is $\rho(q_0)=F+\delta_0$, where $F$ is the frozen part and $\delta_0$ the initial
adapter, and the defect is $F+\delta_0-W$.
\begin{itemize}
  \item \emph{Split first.} $F=\kappa(W_{\rm res})$ with $W_{\rm res}=W-\delta_0$, so the defect is
  $\kappa(W_{\rm res})+\delta_0-W=\kappa(W_{\rm res})-W_{\rm res}$. In LoftQ's quantisation half-step for $t\ge2$ the same
  computation applies with $\delta_0$ replaced by the current fit $\delta(q_{t-1})=L_{t-1}R_{t-1}$.
  \item \emph{Quantise, then fit.} $F=\kappa W$ and HF sets $L R=\operatorname{SVD}_r(W-\kappa W)=-\operatorname{SVD}_r(e)$
  (the best rank-$r$ approximation is odd, $\operatorname{SVD}_r(-X)=-\operatorname{SVD}_r(X)$), and the forward pass
  multiplies by $s$, so $\delta_0=-s\operatorname{SVD}_r(e)$ and the defect is $e-s\operatorname{SVD}_r(e)$. Write
  $P=\operatorname{SVD}_r(e)$. By \auxref{B.11}(b), $P$ and $e-P$ are Frobenius-orthogonal, so
  $\lVert e-sP\rVert^2=\lVert(e-P)+(1-s)P\rVert^2=\lVert e-P\rVert^2+(1-s)^2\lVert P\rVert^2$. For $e\ne0$, $P\ne0$, so the
  defect equals $e-P$ iff $s=1$; then it has rank $\rk e-r$ when $\rk e\ge r$, being the sum of the singular terms beyond
  the $r$-th. HF's defaults $r=\alpha=8$ give $s=\alpha/r=1$. For $s=2$, $\lVert e-2P\rVert^2=\lVert e-P\rVert^2+\lVert
  P\rVert^2=\lVert e\rVert^2$, again by orthogonality; in general $\lVert e-sP\rVert<\lVert e\rVert$ iff $(1-s)^2<1$, that is
  $0<s<2$. A module rolled back to the zero initialisation has $\delta_0=0$ and defect $\kappa W-W=e$.
  \item \emph{Quantise, split, re-quantise.} The split is computed from $\kappa W$, the residual $r'=\kappa W-\delta_0$ is
  re-quantised, and $F=\kappa(r')$. The defect is $\kappa(r')+\delta_0-W=(\kappa(r')-r')+(\kappa W-W)=e+\kappa(r')-r'$.
\end{itemize}
For quantise-then-split with frozen part $F$, the defect $F+\delta_0-W$ equals $e=\kappa W-W$ iff $F=\kappa W-\delta_0$.
Storing $\kappa W$ quantised together with the frozen factors of $\delta_0=sL_0R_0$, which hold $r(m+n)$ numbers, realises
this $F$ exactly. Re-quantising instead adds $\kappa(r')-r'$, which vanishes iff $r'$ is representable, that is, iff $r'$
lies in the null set of representable matrices.

\emph{Action families.} OFT at $R=I$ gives $\rho(q_0)=\kappa\theta_0\,I=\kappa\theta_0$, BOFT at $R=I$ with unit scale gives
$\diag(\mathbf 1)\,\kappa\theta_0\,I=\kappa\theta_0$, and HiRA at $B=0$ gives $\kappa\theta_0+\kappa\theta_0\odot(0\cdot
A)=\kappa\theta_0$; each has defect $e$. DoRA at $B=0$ computes $\rho(q_0)=\diag\big(\mu_{0,i}/\lVert(\kappa\theta_0)_i\rVert\big)
\kappa\theta_0$, normalising the rows of its frozen base $\kappa\theta_0$ (assumed nonzero) and rescaling them by the
magnitude vector $\mu_0$. If $\mu_0$ holds the row norms of $\kappa\theta_0$, then $\rho(q_0)=\kappa\theta_0$ and the defect
is $e$; if it holds those of $\theta_0$, the defect is the displayed matrix.

\emph{Moved invariants.} HiRA's update $W\mapsto W\odot(\mathbf 1+BA)$ preserves every zero entry of its base
(\thmref{Proposition}{II.9}(e)). NF4 contains the code $0$, so every entry of $\theta_0$ whose ratio to its block's absmax
rounds to $0$ becomes an exact zero of $\kappa\theta_0$, while zero entries of $\theta_0$ stay zero; the zero set of
$\kappa\theta_0$ therefore contains that of $\theta_0$, and generically more. Exact-Cayley OFT has
$\rho(Q)\rho(Q)\T=\kappa\theta_0\,C(Q)C(Q)\T\kappa\theta_0\T=\kappa\theta_0\kappa\theta_0\T$ (\thmref{Corollary}{II.6}). HF's
default factor is $R=C(Q)(I-Q^4)$ (\auxref{B.23}(e)), so $\rho\rho\T=\kappa\theta_0(I-Q^4)^2\kappa\theta_0\T=
\kappa\theta_0\kappa\theta_0\T+O(\lVert Q\rVert^4)$.

(c) LoftQ starts from $L_0R_0=0$ and, for $t\ge1$, sets $Q_t=\kappa(\theta_0-L_{t-1}R_{t-1})$ and
$L_tR_t=\operatorname{SVD}_r(\theta_0-Q_t)$. By \auxref{B.11}(b) the SVD step minimises $\lVert\theta_0-Q_t-LR\rVert$ over
rank-$r$ products, so it never increases the objective. The quantisation step returns $\kappa(y)$ for $y=\theta_0-L_{t-1}
R_{t-1}$, which need not be the representable matrix nearest to $y$, and the objective can increase. An example with blocks
equal to the rows of a $2\times3$ matrix, absmax scales, the codebook $\{0,\pm\tfrac12,\pm1\}$ and $r=1$:
\[\theta_0=\begin{pmatrix}1&1&-\tfrac54\\ \tfrac54&1&-1\end{pmatrix},\quad
Q_1=\kappa(\theta_0)=\tfrac54\begin{pmatrix}1&1&-1\\1&1&-1\end{pmatrix},\quad
\theta_0-Q_1=\tfrac14\begin{pmatrix}-1&-1&0\\0&-1&1\end{pmatrix}.\]
Here $(\theta_0-Q_1)(\theta_0-Q_1)\T=\tfrac1{16}\left(\begin{smallmatrix}2&1\\1&2\end{smallmatrix}\right)$ has eigenvalues
$\tfrac3{16}$ and $\tfrac1{16}$, with top eigenvector $(1,1)/\sqrt2$, so $L_1R_1=\tfrac12\left(\begin{smallmatrix}1&1\\1&1
\end{smallmatrix}\right)(\theta_0-Q_1)=\tfrac18\left(\begin{smallmatrix}-1&-2&1\\-1&-2&1\end{smallmatrix}\right)$ and the
objective after the SVD step is $\lVert\theta_0-Q_1-L_1R_1\rVert^2=\tfrac1{16}$. The next quantisation step acts on
$y=\theta_0-L_1R_1=\tfrac18\left(\begin{smallmatrix}9&10&-11\\11&10&-9\end{smallmatrix}\right)$. Each row has absmax
$\tfrac{11}8$, and the ratios $\tfrac9{11}$ and $\tfrac{10}{11}$ round to $1$, so $Q_2=\tfrac{11}8\left(\begin{smallmatrix}1&1&-1\\1&1&-1
\end{smallmatrix}\right)$ and
\[\lVert\theta_0-Q_2-L_1R_1\rVert^2=\lVert y-Q_2\rVert^2=2\Big(\tfrac1{16}+\tfrac1{64}\Big)=\tfrac5{32}>\tfrac1{16}.\]
The objective rose, and $Q_1$, which is representable, is closer to $y$ than $\kappa(y)$. The deployed weight after $T$
iterations is $Q_T+sL_TR_T$, because HF stores the fitted factors and the forward pass multiplies by $s$; its defect is
\[Q_T+sL_TR_T-W=(Q_T+L_TR_T-W)+(s-1)L_TR_T .\]
For $s=1$ this is minus the attained residual, so the output is pointed up to that residual. For $s\ne1$ the second term
$(s-1)L_TR_T$ is proportional to the fit itself and does not shrink with the residual; in the referee simulation of the official configuration ($r=64$, $\alpha=16$, $s=1/4$,
five iterations, a Gaussian $512\times512$ NF4 weight) the attained residual falls to $0.67\lVert e\rVert$ while the deployed
defect rises from $0.91\lVert e\rVert$ to $0.96\lVert e\rVert$ (Appendix~\ref{app:repro}).

(d) Fine-tuning and then quantising gives $\kappa(\theta_0+\Delta)$, a value of $\kappa$. Exact merging after quantising gives
$\kappa\theta_0+\Delta$, which is representable iff $\Delta$ lies in the translate of the representable set by
$-\kappa\theta_0$, a null set; so for generic $\Delta$ it leaves the quantised type and the two orders differ. For a uniform
grid with a fixed scale $a$ per block, the representable matrices are $a$ times integer codes in range, so
$\kappa\theta_0+\Delta$ is representable iff $\Delta=a(c'-c_0)$ with $c'$ in range, that is, iff $\Delta$ lies on the scale
lattice. NF4 is not closed under addition. Its smallest positive code is $q_1\approx0.0796$, and $2q_1\approx0.1592$ is not a
code (the next codes are $0.0796$ and $0.1609$). The blocks $(1,q_1,0,\dots)$ and $(0,q_1,0,\dots)$ are representable,
with absmax scales $1$ and $q_1$, but their sum $(1,2q_1,0,\dots)$ has absmax $1$ and needs the code $2q_1$. Absmax grids
with uniform codes of step $h$ are not closed either: $(1,h,0,\dots)$ and $(h,0,0,\dots)$ are representable, while their sum
$(1+h,h,0,\dots)$ would need the code $h/(1+h)$, which is not a multiple of $h$ for $0<h<1$. A type-preserving merge
returns $\kappa(\kappa\theta_0+\Delta)$, which is representable and differs from $\kappa\theta_0+\Delta$ whenever the latter
is not, so it is lossy for generic $\Delta$. Finally, affine group quantisation stores each group $g$ as $a_gc_g+\beta_g\mathbf
1$, with integer codes $c_g$ and a high-precision scale $a_g$ and zero-point $\beta_g$. QA-LoRA pools its input over the
quantisation groups, so its update is constant on each group, $\Delta_g=\delta_g\mathbf 1$ \citep{xu2023qa}, and
$a_gc_g+\beta_g\mathbf 1+\delta_g\mathbf 1=a_gc_g+(\beta_g+\delta_g)\mathbf 1$ keeps the codes and moves only the zero-point.
\end{proof}

% ====================================================================================================
\subsection{Theorem V.3}\label{proof:V.3}

\begin{restate}{Theorem}{V.3}{rebasing closure}
(a) \emph{Additive methods with a fixed shape $C\ni0$}, with frames, masks and compression operators frozen at $\theta_0$
across cycles. Then $R_K=\theta_0+C^{+K}$, the $K$-fold Minkowski sum, and $R_\infty=\theta_0+\operatorname{Mon}(C)$, the
additive monoid generated by $C$. If $C$ is closed under real scalings, $R_\infty=\theta_0+\spanop C$.

(b) \emph{Action methods}, where one cycle applies an element of a set $S$ with $1\in S\subseteq G$; for scaling methods,
$S\subseteq\operatorname{End}(\Theta)$ is a subset containing the identity. Then $R_K=S^K\theta_0$, and $R_\infty$ is the
orbit of the monoid generated by $S$. Suppose $S$ lies in the subgroup generated by connected Lie subgroups
$H_1,\dots,H_p$ and contains a neighbourhood of $1$ in each. Then the monoid generated by $S$ is the connected Lie subgroup
$H$ with $\operatorname{Lie}(H)=\operatorname{Lie}\langle\mathfrak h_1,\dots,\mathfrak h_p\rangle$.

(c) \emph{Block-orthogonal factors.} Let the factors have coordinate blocks $I_j\subseteq\{1,\dots,n\}$ (block OFT, BOFT
factors, Givens planes). Then $\operatorname{Lie}\langle\so(I_j)\rangle_j=\bigoplus_C\so(C)$, summed over the connected
components $C$ of the hypergraph $(\{1,\dots,n\},\{I_j\})$.

\emph{Corollaries.}
\begin{enumerate}[label=\arabic*.,leftmargin=1.6em]
  \item \textbf{ReLoRA.} $R_K=W_0+\MM_{\le\min(Kr,m,n)}$. It reaches full fine-tuning iff $K\ge\lceil\min(m,n)/r\rceil$.
  \item \textbf{Idempotent methods} gain nothing from rebasing. These are the methods whose shape is exactly a subspace,
  subgroup or submonoid independent of the base, with frames, masks and compression operators frozen across cycles:
  LoRA-FA with a fixed $A$, LoRA-XS with frozen frames, MoRA with fixed operators, FourierFT with fixed indices, WaveFT,
  C3A, MiSS (default balance mode; bat mode is base-dependent but still idempotent, since its blockwise maps
  $W\mapsto(W+I)(I+M)-I$ form a monoid), SHiRA, BitFit, FISH with a fixed mask, LN-tuning, (IA)$^3$, SSF and RoAd. Block OFT
  with a fixed partition and an exact Cayley map gains only the measure-zero set of block rotations with an eigenvalue
  $-1$, so $R_\infty=W_0\cdot\prod\SO(b)=\overline{R_1}$. Constrained OFT (COFT, an $\varepsilon$-ball) with an exact Cayley
  map gains up to $\prod\SO(b)$; HF's \texttt{coft=True} keeps the Cayley--Neumann map by default and so falls under
  corollary~4. Methods that recompute frames or masks from the merged weight (an SVD-recomputing LoRA-XS, ReMoRA's
  alternating operators, resampled indices) do gain.
  \item \textbf{VeRA.} $\mathrm{VeRA}^\infty=\text{LoRA-FA}(A)$ whenever the frozen $B$ has no zero entries.
  \item \textbf{OFT family.} Block OFT with an exact Cayley map is stuck in $W_0\cdot\prod\SO(b)$. HF's default
  Cayley--Neumann factor is $R=C(Q)(I-Q^4)$, with $C$ the exact Cayley map, so $R\T R=(I-Q^4)^2$ and $R$ is not orthogonal.
  Its singular values are $\lvert1-\mu_k^4\rvert$, where $\pm i\mu_k$ are the eigenvalues of $Q$. In the operating regime
  $\lVert Q\rVert_{\rm op}\le2^{1/4}$, every factor is a contraction. Re-merging then decreases $K_{\rm out}$ in the
  Loewner order and no singular value grows, so the drift is monotone shrinkage, and the reachable set lies in $W_0$ times
  the block contraction semigroup. With unconstrained $Q$ the factors include singular ones ($\mu_k=1$), and re-merging
  generates $\prod\big(\GL^+(b)\cup\MM_{\le b-2}\big)$, or $\prod\big(\CO^+(2)\cup\{0\}\big)$ when $b=2$; this clause is
  established at the level of a proof sketch. Iterated BOFT (exact Cayley in HF) with $m'$ butterfly factors
  of block size $b$ reaches $\diag(\R^m)\cdot W_0\cdot P\big(\prod\SO(b\,2^{m'-1})\big)P\T$, because the factor hypergraph
  has components of size $b\,2^{m'-1}$. For injective $W_0$ this is $\diag(\R^m)\cdot W_0\cdot\SO(n)$ iff
  $b\,2^{m'-1}=n$, the full-depth butterfly $m'=1+\log_2(n/b)$. HF's default \texttt{boft\_\allowbreak n\_\allowbreak butterfly\_\allowbreak factor=1} is
  block OFT times $\diag(s)$. A \emph{single} full-depth BOFT has dimension at most $m'\tfrac nb\binom b2$, which is below
  $\dim\SO(n)$ when $m'(b-1)<n-1$; its matrices are dense, while the set they form has positive codimension.
  \item \textbf{Re-HRA} (HF \texttt{apply\_GS=False}), with $W_0$ injective and $r$ even. Then
  $R_K=(W_0+\MM_{\le Kr})\cap W_0\cdot\SO(n)$, so rebasing commutes with the meet of \thmref{Theorem}{II.7}. This
  multiplicative twin of ReLoRA reaches $W_0\cdot\SO(n)$ in at most $\lceil n/r\rceil$ cycles.
  \item \textbf{GaLore (with a gradient-SVD or a random projector), LISA and MeZO}, under the optimizer hypotheses of
  \thmref{Proposition}{IV.2}, are rebasing schemes of \emph{linear} methods whose shape moves. That is why, as a set over
  the sequences of shapes they can choose, they reach all of $\Theta$, while a fixed linear shape $L$ reaches only
  $\theta_0+L$ (corollary~2). Low rank is the per-period obstruction only for GaLore; LISA and MeZO already make full-rank
  per-matrix updates and gain coverage. Flora proper takes full-space steps with a compressed momentum, so it has no
  per-period low-rank shape.
  \item \textbf{Merging never rescues an exactly orthogonal method from its absolute invariants.} In exact arithmetic,
  every exactly orthogonal input-side rebasing preserves $K_{\rm out}$ and the spectrum, and every output-side one
  preserves $K_{\rm in}$ and the spectrum.
\end{enumerate}
\end{restate}

\begin{proof}
(a) With frames, masks and compression operators frozen, every cycle trains the same additive shape over the current
base: $\Im M(\theta)=\theta+C$ for every $\theta$. By the recursion of \S\ref{proof:V.1}, $R_1=\theta_0+C$ and
$R_{K+1}=\bigcup_{\theta\in R_K}(\theta+C)=R_K+C$, so $R_K=\theta_0+C^{+K}$ by induction. Since $0\in C$, the sets $C^{+K}$
increase, and their union is the additive monoid $\operatorname{Mon}(C)$ generated by $C$ (\bgref{D.10}); hence
$R_\infty=\theta_0+\operatorname{Mon}(C)$. If $\lambda C\subseteq C$ for every real $\lambda$, then $-C=C$, so
$\operatorname{Mon}(C)$ is closed under negation and is an additive subgroup, and
$\lambda(x_1+\dots+x_K)=\lambda x_1+\dots+\lambda x_K$ shows that it is closed under scaling. So it is a linear subspace
containing $C$ and contained in $\spanop C$, and equals $\spanop C$.

(b) Each cycle applies one element of $S$ to the current point, so $\Im M(\theta)=S\theta$, and the recursion gives
$R_K=S^K\theta_0$. Since $1\in S$, the sets $S^K$ increase and their union is the monoid $\operatorname{Mon}(S)$ generated
by $S$, so $R_\infty=\operatorname{Mon}(S)\,\theta_0$, the orbit of $\theta_0$ under that monoid. For the second statement,
let $U_i\subseteq S$ be a neighbourhood of $1$ in $H_i$. Then $U_i\cap U_i^{-1}$ is a symmetric neighbourhood of $1$ in the
connected Lie group $H_i$, so it generates $H_i$ as a monoid (\auxref{B.24}(a)), and $H_i\subseteq\operatorname{Mon}(S)$.
Each $H_i$ is a group, so the monoid generated by $H_1\cup\dots\cup H_p$ is closed under inverses and equals the subgroup
they generate, which is the connected Lie subgroup $H$ with $\operatorname{Lie}(H)=\operatorname{Lie}\langle\mathfrak
h_1,\dots,\mathfrak h_p\rangle$ (\auxref{B.24}(d)). So $H\subseteq\operatorname{Mon}(S)$. Conversely $S\subseteq H$ by
hypothesis, and $H$ is closed under products, so $\operatorname{Mon}(S)\subseteq H$.

(c) Write $L_{ab}=E_{ab}-E_{ba}$, so $\so(I)=\spanop\{L_{ab}:a,b\in I\}$. For distinct $a,b,c$,
\[L_{ac}L_{cb}=E_{ab},\qquad L_{cb}L_{ac}=E_{ba},\qquad\text{so}\qquad[L_{ac},L_{cb}]=L_{ab},\]
while $L_{ab}L_{cd}=L_{cd}L_{ab}=0$ when $\{a,b\}\cap\{c,d\}=\emptyset$. Let $\mathfrak g=\operatorname{Lie}\langle\so(I_j)
\rangle_j$. If $a\ne b$ lie in one component $C$, join them by a shortest chain $a=c_0,c_1,\dots,c_t=b$ of distinct
coordinates with consecutive ones in a common block. By induction on $t$, $L_{ab}\in\mathfrak g$: for $t=1$,
$L_{ab}\in\so(I_j)$ for a block containing both; for $t>1$, $L_{ac_{t-1}}\in\mathfrak g$ by induction and $L_{c_{t-1}b}\in
\mathfrak g$, and their bracket is $L_{ab}$. So $\bigoplus_C\so(C)\subseteq\mathfrak g$. Conversely, $\bigoplus_C\so(C)$ is a
Lie subalgebra, since each $\so(C)$ is one and elements of different components have disjoint supports and commute, and it
contains every $\so(I_j)$, because each block lies in one component. Hence $\mathfrak g=\bigoplus_C\so(C)$, of dimension
$\sum_C\binom{\lvert C\rvert}2$.

\emph{Corollary 1.} LoRA's shape $\{sBA\}$ is $\Mr$, which contains $0$ and is closed under real scalings, so (a) applies.
By \auxref{B.12}(a), a sum of $K$ matrices of rank at most $r$ has rank at most $\min(Kr,m,n)$. Conversely, a matrix of rank
$k\le\min(Kr,m,n)$ is a sum of $k$ rank-one terms of its SVD (\auxref{B.11}(a)), which split into $K$ groups of at most $r$
terms, each in $\Mr$. So $C^{+K}=\MM_{\le\min(Kr,m,n)}$. This is all of $\R^{m\times n}$ iff $Kr\ge\min(m,n)$, that is,
iff $K\ge\lceil\min(m,n)/r\rceil$.

\emph{Corollary 2.} If the shape is an additive submonoid $C$ independent of the base (in particular a linear subspace),
then $C^{+K}=C$ and $R_K=R_1$ for all $K$ by (a); if it is a submonoid $S$ of a group or of $\operatorname{End}(\Theta)$,
then $S^K=S$ and $R_K=R_1$ by (b). The listed methods are of this kind once their frames, masks and operators are frozen:
the updates of LoRA-FA ($\{XA\}$), LoRA-XS ($\{UMV\T\}$), MoRA (linear in its square matrix), FourierFT and WaveFT (fixed
coefficient positions), C3A (block-circulant matrices), MiSS in balance mode (linear in its shared block), SHiRA, BitFit,
FISH and LN-tuning (fixed coordinates) form linear subspaces; the scalings of (IA)$^3$, $W\mapsto\diag(\ell)W$, form a
monoid in $\operatorname{End}(\Theta)$; SSF's maps $y\mapsto s\odot y+t$ form a monoid under composition,
$s'\odot(s\odot y+t)+t'=(s's)\odot y+(s't+t')$; RoAd's $2\times2$ blocks range over the scaled rotations or over all
$2\times2$ matrices, both monoids; and MiSS's bat-mode maps satisfy $(W+I)(I+M)(I+M')-I$ with $(I+M)(I+M')=I+M''$,
$M''=M+M'+MM'$ again in the class, so they form a monoid.

For block OFT with an exact Cayley map and a fixed partition, $S=\prod\mathrm{Cay}(b)$ with
$\mathrm{Cay}(b)=C(\so(b))=\{R\in\SO(b):-1\notin\operatorname{spec}R\}$, open and dense in $\SO(b)$ (\auxref{B.23}(c)). Every
element of $\SO(b)$ is a product of two elements of $\mathrm{Cay}(b)$ (\auxref{B.23}(d)), so $S^2=\prod\SO(b)$, and
$R_K=W_0\cdot\prod\SO(b)$ for $K\ge2$. By continuity of $g\mapsto W_0g$ and compactness of $\prod\SO(b)$,
$\overline{R_1}=W_0\cdot\overline{\prod\mathrm{Cay}(b)}=W_0\cdot\prod\SO(b)$, and $R_\infty\setminus R_1$ lies in the image
of the Haar-null set $\{\prod_j\det(I+R_j)=0\}$. For COFT with an exact Cayley map, $S$ is the image under $C$ of a ball
about $0$, which contains a neighbourhood of $1$ in $\prod\SO(b)$ because $C$ is a local diffeomorphism at $0$; by (b) with
the single connected group $\prod\SO(b)$, $\operatorname{Mon}(S)=\prod\SO(b)$.

Methods whose shape changes with the base escape (a). For linear shapes $L_1\ne L_2$ used in two cycles (resampled indices,
alternating operators), $R_2\supseteq\theta_0+L_1+L_2\supsetneq\theta_0+L_1=R_1$ whenever $L_2\not\subseteq L_1$. For LoRA-XS
with frames recomputed from the merged weight, take $r=1$ and $W_0=\diag(2,1)$: $R_1=\{\diag(2+t,1)\}$; the point
$\diag(-\tfrac12,1)\in R_1$ has top singular direction $e_2$, so the second cycle adds $\{te_2e_2\T\}$ and reaches
$\diag(-\tfrac12,1+t)\notin R_1$.

\emph{Corollary 3.} VeRA's shape is $C=\{\Lambda_bB\Lambda_dA\}$ with $B$, $A$ frozen and $b\in\R^m$, $d\in\R^r$ trainable. It
contains $0$, is closed under real scalings ($b\mapsto\lambda b$), and lies in $\{XA:X\in\R^{m\times r}\}$. The entries of
$\Lambda_bB\Lambda_d$ are $b_iB_{ij}d_j$, so $b=e_i$, $d=e_j$ gives $B_{ij}E_{ij}$; if every $B_{ij}\ne0$, the span of the
matrices $\Lambda_bB\Lambda_d$ is $\R^{m\times r}$, and $\spanop C=\{XA\}$. By (a), $R_\infty=W_0+\{XA\}$, the image of
LoRA-FA with the same $A$.

\emph{Corollary 4.} \emph{Exact Cayley.} This is corollary~2.

\emph{The Cayley--Neumann factor.} By \auxref{B.23}(e), $R=C(Q)(I-Q^4)$, $R\T R=(I-Q^4)^2$, and the singular values of $R$ are
$\lvert1-\mu_k^4\rvert$, each twice, and $1$ for a zero eigenvalue of $Q$. So $R$ is orthogonal only when every $\mu_k^4\in
\{0,2\}$. If $\lVert Q\rVert_{\rm op}\le2^{1/4}$, then $0\le\mu_k^4\le2$ and $\lvert1-\mu_k^4\rvert\le1$, so $R$ is a contraction.
For an input-side merge $W\mapsto WR$ with block-diagonal $R$,
\[K_{\rm out}(WR)=WRR\T W\T=W(I-Q^4)^2W\T\preceq WW\T=K_{\rm out}(W),\]
because $C(Q)$ commutes with $Q^4$ and $(I-Q^4)^2\preceq I$ (\auxref{B.18}(b)); by \auxref{B.18}(c) no singular value of $W$
grows. Products of block-diagonal contractions are block-diagonal contractions (\auxref{B.18}(a)), so the reachable set lies
in $W_0$ times that semigroup.

\emph{Unconstrained $Q$.} Fix one block of size $b$ and let $S_b=\{R(Q):Q\in\so(b)\}$; since the blocks are independent and
$R(0)=I$, the monoid generated by the block-diagonal factors is $\prod\operatorname{Mon}(S_b)$. Write $Q=O(\bigoplus_k\mu_kJ
\oplus0)O\T$ with $O\in\Ort(b)$, $J=\left(\begin{smallmatrix}0&-1\\1&0\end{smallmatrix}\right)$ and $\mu_k\ge0$; then
$I-Q^4=O(\bigoplus_k(1-\mu_k^4)I_2\oplus I)O\T$.

\emph{$\subseteq$.} $\det R=\det C(Q)\det(I-Q^4)=\prod_k(1-\mu_k^4)^2\ge0$, and $\rk R=\rk(I-Q^4)=b-2\#\{k:\mu_k=1\}$. So an
invertible factor lies in $\GL^+(b)$ and a singular one has rank at most $b-2$. Products of elements of $\GL^+(b)$ stay in
it, and any product with a factor of rank at most $b-2$ has rank at most $b-2$. Hence $\operatorname{Mon}(S_b)\subseteq\GL^+
(b)\cup\MM_{\le b-2}$.

\emph{Positive scalings.} $C(-Q)=C(Q)^{-1}$ commutes with $Q^4$, so $R(Q)R(-Q)=(I-Q^4)^2=O(\bigoplus_k(1-\mu_k^4)^2I_2\oplus
I)O\T$. As $\mu_k$ ranges over $[0,\infty)$, $(1-\mu_k^4)^2$ takes every positive value ($\mu_k^4=1+\sqrt\lambda$ gives
$\lambda$). Taking $O$ a permutation matrix, $\operatorname{Mon}(S_b)$ contains every positive diagonal matrix that is
constant on the pairs of a matching of the coordinates and $1$ off it. For $b\ge3$ these generate all positive diagonal
matrices: their logarithms include every multiple of $e_i+e_j$, $i\ne j$, and $e_i=\tfrac12\big((e_i+e_j)+(e_i+e_k)-(e_j+
e_k)\big)$, so products of commuting such matrices realise every $\diag(e^x)$. For $b=2$ they are the positive scalars.

\emph{Rotations.} For $\lVert Q\rVert_{\rm op}<1$, $I-Q^4$ is positive definite with eigenvalues $1-\mu_k^4$ on the planes
of $Q$. Choose $Q'=O(\bigoplus_k\mu_k'J\oplus0)O\T$ with the same $O$ and $(1-\mu_k'^4)^2=(1-\mu_k^4)^{-1}$; then
$R(Q')R(-Q')=(I-Q^4)^{-1}$, and $C(Q)=R(Q)R(Q')R(-Q')\in\operatorname{Mon}(S_b)$. The set $\{C(Q):\lVert Q\rVert_{\rm op}<1\}$
is a neighbourhood of $I$ in $\SO(b)$ (\auxref{B.23}(c)), symmetric because $C(-Q)=C(Q)^{-1}$, so it generates $\SO(b)$ as a
monoid (\auxref{B.24}(a)). With the positive diagonals and the decomposition $\GL^+(b)=\SO(b)\,D^+\SO(b)$
(\auxref{B.19}(c)), $\GL^+(b)\subseteq\operatorname{Mon}(S_b)$ for $b\ge3$; for $b=2$, $\SO(2)\cdot\R_{>0}=\CO^+(2)$.

\emph{Singular blocks.} Let $V\subseteq\R^b$ have dimension $b-2$ and choose $O$ mapping the last coordinate plane onto
$V^\perp$, and $Q_V=O(0\oplus J)O\T$. Then $Q_V^4$ is the projector onto $V^\perp$, $C(Q_V)$ acts as $J$ on $V^\perp$ (since
$(I+J)(I-J)^{-1}=J$) and as the identity on $V$, and $R(Q_V)=C(Q_V)(I-Q_V^4)=P_V$, the orthogonal projector onto $V$. If
$P_U$ with $\dim U=k$, $1\le k\le b-2$, is in the monoid, choose $U'\subseteq U$ of dimension $k-1$ and $W'\subseteq U^\perp$
of dimension $b-1-k\le\dim U^\perp$, and put $V=U'\oplus W'$, of dimension $b-2$. Since $W'\perp U$,
$P_UP_V=P_U(P_{U'}+P_{W'})=P_{U'}$. By downward induction every orthogonal projector of rank at most $b-2$ lies in the
monoid. A matrix $X$ of rank $k\le b-2$ is $X=U\diag(\sigma_1,\dots,\sigma_k,1,\dots,1)\,P_k\,V\T$ by its SVD, with $P_k$
the projector onto the first $k$ coordinates. Changing the sign of the last column of $U$, or of the last row of $V\T$,
leaves this product unchanged, because $P_k$ kills the last coordinate; so both outer factors may be taken in $\GL^+(b)$,
and $X\in\operatorname{Mon}(S_b)$. For $b=2$, $R(\mu J)=(1-\mu^4)\,C(\mu J)$ is a real multiple of a rotation, so
$\operatorname{Mon}(S_2)\subseteq\CO^+(2)\cup\{0\}$, and $R(J)=0$. Altogether
$\operatorname{Mon}(S_b)=\GL^+(b)\cup\MM_{\le b-2}$ for $b\ge3$ and $\CO^+(2)\cup\{0\}$ for $b=2$, which is the clause; the
steps above complete the sketch recorded in the main text.

\emph{BOFT.} Group the coordinates into runs of $b/2$ consecutive coordinates, indexed by $p=0,1,\dots,2n/b-1$. Up to one fixed
relabelling of coordinates, the $l$-th butterfly factor of BOFT ($l=1,\dots,m'$) is block-diagonal with one orthogonal
$b\times b$ block, an exact Cayley image, on each pair of runs whose indices differ exactly in binary digit $l-1$
\citep{liu2023parameter,dao2019learning}; for $l=1$ these are $n/b$ consecutive blocks of size $b$, block OFT. Two runs lie
in one component of the factor hypergraph iff their indices agree in all digits from $m'$ on, so the components consist of
$2^{m'}$ runs, $b\,2^{m'-1}$ coordinates. HF's BOFT multiplies the output by a trainable scale, $W\mapsto\diag(s)\,W\,R$, and
left and right multiplications commute, so after $K$ cycles the weight is $\diag(s_K\cdots s_1)\,W_0\,R_1\cdots R_K$; the
products of scale vectors already fill $\R^m$ at $K=1$. The orthogonal part is generated by the set $S$ of products of the
$m'$ butterfly factors. Let $H_l$ be the group of all block rotations of level $l$, connected; $S$ lies in the group they
generate, and setting all but one factor to the identity shows that $S$ contains a neighbourhood of $1$ in each $H_l$, the
exact Cayley image being open. By (b) and (c), $\operatorname{Mon}(S)$ is the connected Lie subgroup with Lie algebra
$\bigoplus_C\so(C)$, which is the group of block rotations on the components, $P\big(\prod\SO(b\,2^{m'-1})\big)P\T$ for the
permutation $P$ that makes the components consecutive. This proves the reach. If $b\,2^{m'-1}=n$ there is one component and
the reach is $\diag(\R^m)\cdot W_0\cdot\SO(n)$. If $c:=b\,2^{m'-1}<n$, the group $G=P\big(\prod\SO(c)\big)P\T$ has dimension
$n(c-1)/2<n(n-1)/2$; and for $W_0$ of full column rank the reachable sets differ (for $W_0=0$ both are $\{0\}$). Indeed, if
$\diag(\R^m)W_0G=\diag(\R^m)W_0\SO(n)$, then each $R\in\SO(n)$ satisfies $W_0R=\diag(s)W_0g$ for some $s$ and $g\in G$, so
$M=Rg^{-1}$ satisfies $W_0M=\diag(s)W_0$: every row $w_i$ of $W_0$ is an eigenvector of $M\T$. These rows span $\R^n$, so $M\T$ is
diagonalisable with real eigenvalues; being orthogonal, $M$ is an orthogonal involution $I-2P_E$ whose $-1$-eigenspace $E$ is
spanned by the rows of $W_0$ it contains. So $M$ ranges over a finite set $F_0$, and $\SO(n)\subseteq F_0G$, a finite union
of submanifolds of lower dimension, which is impossible. With $m'=1$ the single level is block OFT, so HF's default
\texttt{boft\_n\_butterfly\_factor=1} is block OFT times $\diag(s)$.

\emph{A single BOFT.} Its parameters number $m'\frac nb\binom b2$, so its image has at most this dimension
(\auxref{B.9}(b)), and $m'\frac nb\frac{b(b-1)}2<\frac{n(n-1)}2$ iff $m'(b-1)<n-1$; the image then has positive codimension
in $\SO(n)$ and is Haar-null. At full depth, $2^{m'}=2n/b$ runs form one component, and for every pair of coordinates $(i,j)$
there is a chain $j=c_0,c_1,\dots,c_{m'}=i$ in which $c_{l-1}$ and $c_l$ lie in a common block of the $l$-th factor applied:
each factor lets the run index flip one binary digit, a different one for each factor, and lets the position inside the run
be chosen freely. So the entry $(i,j)$ of the product, a polynomial in the block entries, is nonzero at a product of signed
permutation matrices of determinant $1$ (one per block) that routes $j$ to $i$ along the chain. It is therefore not
identically zero on the connected group $\prod_lH_l$, and since the Cayley parametrisation is a diffeomorphism onto an open
subset of that group, the entry is nonzero for parameters outside a nowhere-dense analytic subset (\auxref{B.10}). All
$n^2$ entries are nonzero for generic parameters, so the matrices are dense.

\emph{Corollary 5.} With \texttt{apply\_GS=False}, one HRA cycle applies $H_{u_1}\cdots H_{u_r}$ on the input side. Let
$S_k=\{R\in\SO(n):\rk(R-I)\le k\}$ for even $k$. A product of $r$ reflections lies in $\SO(n)$ ($r$ even) and has
$\rk(R-I)\le r$ (\auxref{B.22}); conversely, if $R\in S_r$ then $j=\rk(R-I)$ is even and $R$ is a product of exactly $j$
reflections (\auxref{B.22}), which pairs $H_uH_u=I$ pad to $r$. So one cycle applies exactly the elements of $S_r$, and by
(b) $R_K=W_0\cdot S_r^K$. Next, $S_r^K=S_{Kr}$. For $\subseteq$, $R_1R_2-I=R_1(R_2-I)+(R_1-I)$ gives
$\rk(R_1R_2-I)\le\rk(R_2-I)+\rk(R_1-I)$ (\auxref{B.12}(a)). For $\supseteq$, write $R\in S_{Kr}$ as a product of
$j=\rk(R-I)\le Kr$ reflections, $j$ even, split them into $K$ consecutive groups of even size at most $r$, and pad each group
with pairs to size $r$. Finally, since $W_0$ is injective, $\rk(W_0R-W_0)=\rk(R-I)$, so
$W_0\cdot S_{Kr}=(W_0+\MM_{\le Kr})\cap W_0\cdot\SO(n)$. When $Kr\ge n$, $S_{Kr}=\SO(n)$, which gives the last sentence.

\emph{Corollary 6.} By \thmref{Proposition}{IV.2} and \S\ref{proof:IV.3}(8), each period of these methods is a period of
the linear method $\theta_k+\mathcal L_kC$ with shape $L_k=\operatorname{im}\mathcal L_k$, so $R_K=\theta_0+L_1+\dots+L_K$
for the chosen shapes. For GaLore with left projection, $L_k=\{Z:\operatorname{col}Z\subseteq\operatorname{col}P_k\}$, and
column spaces that together span $\R^m$ give all of $\R^{m\times n}$ after $\lceil m/r\rceil$ periods; for LISA, layer blocks
covering every layer do the same; for MeZO, $L_k=\R z_k$, and $z_k$ running through a basis give $\Theta$. A fixed linear
shape $L$ is its own additive monoid, so it reaches only $\theta_0+L$ (corollary~2). A GaLore period moves a projected matrix
by an update of rank at most $r$; a LISA period moves each sampled matrix by an unconstrained update, and a MeZO step by a
multiple of the Gaussian $z_k$, whose restriction to each matrix has full rank almost surely. Flora's steps contain a
full-rank term (\S\ref{proof:IV.1}).

\emph{Corollary 7.} For orthogonal $R$, $(WR)(WR)\T=WW\T$ and $(RW)\T(RW)=W\T W$; $WR$ and $RW$ have the singular values of $W$.
By induction over cycles, every input-side rebasing preserves $K_{\rm out}$ and every output-side one $K_{\rm in}$, and both
preserve the spectrum (\thmref{Corollary}{II.6}).
\end{proof}

\begin{remark*}{numerical checks}
The component sizes $4,8,16,32,64$ for $n=64$, $b=4$ and $m'=1,\dots,5$ reported after \thmref{Theorem}{V.3} are the
$b\,2^{m'-1}$ of the BOFT paragraph; the Loewner-monotone decrease of $K_{\rm out}$ over $200$ default-OFT merges, the
contraction bound for random skew $Q$ with $\lVert Q\rVert_{\rm op}\le2^{1/4}$, and $R=0$ at $b=2$, $Q=J$, instantiate the
Cayley--Neumann paragraph (Appendix~\ref{app:repro}).
\end{remark*}

% ====================================================================================================
\subsection{Proposition V.4}\label{proof:V.4}

\begin{restate}{Proposition}{V.4}{combination rules: gauge and naturality}
(a) A rule $\Phi\big((B_i,A_i)_i\big)$ for combining LoRA factors descends to a function of the adapters $(B_iA_i)_i$ iff it
is constant on the fibres of $(B_i,A_i)\mapsto(B_iA_i)$. On the full-rank locus this is equivalent to
$\prod_i\GL_r$-invariance. It is also equivalent for $\Phi$ continuous on the whole factor space; this clause rests on the
unique closed orbit in each fibre, which lies in the closure of every orbit of the fibre.
\begin{itemize}
  \item Task arithmetic $\sum_iw_is_iB_iA_i$, with $s_i$ HF's \texttt{scaling} ($\alpha_i/r_i$, or $\alpha_i/\sqrt{r_i}$
  under rsLoRA), is invariant and lands in $\mathrm{LoRA}_{\sum_ir_i}$. It is HF's \texttt{cat} with merged scale $1$.
  \item LoraHub's factor-wise rule $\big(\sum_iw_iB_i\big)\big(\sum_jw_jA_j\big)$ is not invariant. Replacing $(B_2,A_2)$ by
  $(cB_2,A_2/c)$ changes it by
  \[\sum_{l\ne2}w_lw_2\big[(c-1)B_2A_l+(c^{-1}-1)B_lA_2\big].\]
  HF's \texttt{add\_\allowbreak weighted\_\allowbreak adapter(..., combination\_\allowbreak type="linear")} is a factor-wise
  rule of the same type. It weights $A_i$ by $\operatorname{sign}(w_i)\sqrt{\lvert w_is_i\rvert}$ and $B_i$ by
  $\sqrt{\lvert w_is_i\rvert}$ and sums the factors separately, so it has cross terms. On adapters that share a frozen $A$
  it is still not exact up to a global factor, since the relative weights become $\sqrt{\lvert w_is_i\rvert}$. HF's non-SVD
  \texttt{ties}, \texttt{dare\_*} and \texttt{magnitude\_prune} variants also act on factors; the \texttt{*\_svd} variants
  act on full updates.
  \item Poly and MHR are different. Their skill inventories are trained jointly through the product-of-sums map, which is
  invariant under the diagonal $\GL_r$, one common $g$ for every skill. Numerically, for fixed generic routing with fewer
  skills than tasks and full-rank factors, this is the whole continuous fibre. With at least as many skills as tasks, or
  once the routing is counted as a parameter, there are further symmetries (skill permutations, and re-mixings
  $\alpha\mapsto\alpha M$ with $M\mathbf 1=\mathbf 1$). Their cross terms are a design choice of a jointly trained
  parametrisation, and raise no question of well-definedness.
\end{itemize}

(b) \emph{Well-definedness versus exactness.} A joint gauge fixing makes a factor-wise rule well defined. For
rank-preserving product-of-sums rules, exactness (agreement with task arithmetic for all weights $w$) needs degree one in
$w$, which for injective factors forces a shared factor; \texttt{cat}, which concatenates and raises the rank, is exact
without one. A frozen shared factor (FFA-LoRA; LoRA-FA) gives $\big(\sum_iw_iB_i\big)A=\sum_iw_iB_iA$ exactly. LoraHub with
a shared $A$ still returns $\big(\sum_jw_j\big)\sum_iw_iB_iA$, exact only up to the global factor $\sum_jw_j$, since its
weights are unnormalised. A per-module SVD canonical form leaves a $(\mathbb Z/2)^r$ of sign flips per module, plus
rotations within repeated singular values, and these change the cross terms. What is needed is a \emph{joint} frame that
produces a shared factor, as KnOTS's joint re-factorisation of the concatenated updates through a shared $U\Sigma$ does.
Federated LoRA shows the difference. FedAvg-LoRA clients start each round from the same broadcast $(B,A)$, so their failure,
$\tfrac12(B_1+B_2)\cdot\tfrac12(A_1+A_2)\ne\tfrac12(B_1A_1+B_2A_2)$, is bilinearity, and it persists in a fully fixed gauge.

(c) \emph{Naturality.} On updates with $N$ coordinates, write $T_N$, $B_N$ and $\mathrm{Mon}_N$ for the invertible diagonal,
signed permutation and monomial matrices. Task arithmetic is $\GL(\Theta)$-equivariant on updates. DARE is coordinatewise
and magnitude-free, $D(m\odot\tau)=m\odot(D\tau)$ for every invertible diagonal $D$ and the same mask $m$. So DARE is exactly
equivariant under the torus $T_N$, sign flips included, for each mask, and under permutations in distribution; its
distributional group is exactly $\mathrm{Mon}_N$. TIES with global trimming (top-$k$ magnitude, sign election by total mass
as in the TIES paper and HF's default \texttt{"total"}, disjoint mean) is equivariant under $\R^\times\cdot B_N$,
homotheties times signed permutations, generically (off magnitude ties at the threshold). With HF's
\texttt{majority\_\allowbreak sign\_\allowbreak method="frequency"}, count ties are broken to $+1$ and the group drops to
$\R_{>0}\cdot S_N$. Per-tensor trimming changes the group to $\prod_b\big(\R^\times\cdot B_{N_b}\big)$, extended by
permutations of equal-size blocks. DARE followed by full TIES (with trimming) is $\R^\times\cdot B_N$-equivariant in
distribution; HF's and MergeKit's \texttt{dare\_ties} has no magnitude trim, so with \texttt{"total"} election it is exactly
$T_N$-equivariant per mask, and its group contains $\mathrm{Mon}_N$. None of the sparsifying rules (DARE, TIES, DARE-TIES)
is $\Ort(N)$-equivariant. Applied to factors and summed across adapters (HF's non-SVD variants), all of them are
gauge-dependent, and under independent per-adapter gauges even DARE fails, through the cross terms of (a). DARE's masking
commutes with $T_r$ exactly for each mask and with $\mathrm{Mon}_r$ in distribution, so a single masked adapter, or a
factor-wise sum under one gauge shared by all adapters, is $T_r$-invariant per mask and $\mathrm{Mon}_r$-invariant in
distribution.
\end{restate}

\begin{proof}
Write $\pi\big((B_i,A_i)_i\big)=(B_iA_i)_i$, and let adapter $i$ have rank $r_i$ (the statement's $r$ when all ranks agree);
its gauge group $\GL_{r_i}$ acts by $(B_ig_i^{-1},g_iA_i)$, and $\pi$ is invariant under $\prod_i\GL_{r_i}$.

(a) \emph{Descent.} If $\Phi=\varphi\circ\pi$, then $\Phi$ is constant on the fibres of $\pi$. Conversely, if $\Phi$ is constant
on fibres, $\varphi(\pi(x)):=\Phi(x)$ is a well-defined function on the image of $\pi$.

\emph{Full-rank locus.} Let $\rk B_i=\rk A_i=r_i$ and $B_i'A_i'=B_iA_i$. The product has rank $r_i$, so $\rk B_i'=\rk A_i'=r_i$,
and the full-rank locus is a union of whole fibres. Moreover $\operatorname{col}B_i'=\operatorname{col}(B_iA_i)=
\operatorname{col}B_i$, so $B_i'=B_ig_i^{-1}$ for a unique $g_i\in\GL_{r_i}$; then $B_iA_i=B_ig_i^{-1}A_i'$, and injectivity of
$B_i$ gives $A_i'=g_iA_i$. So each fibre in the full-rank locus is a single $\prod_i\GL_{r_i}$-orbit (as in
\thmref{Theorem}{II.2}), and there constancy on fibres is the same as invariance.

\emph{Continuous rules.} Let $\Phi$ be continuous on the whole factor space and invariant, and let $x,x'$ lie in the fibre
over $(W_i)_i$, $\rk W_i=k_i\le r_i$. By \auxref{B.27}(c), in each factor the pairs with product $W_i$ and both factors of rank
$k_i$ form one $\GL_{r_i}$-orbit $O_{W_i}$, which lies in the closure of every orbit of the fibre. Applying this factor by
factor, there are gauges $g(t),g'(t)\in\prod_i\GL_{r_i}$ with $g(t)\cdot x\to o$ and $g'(t)\cdot x'\to o'$ as $t\to0$,
where $o,o'\in\prod_iO_{W_i}$, and $o'=h\cdot o$ for some gauge $h$. By invariance and continuity,
\[\Phi(x)=\lim_{t\to0}\Phi(g(t)\cdot x)=\Phi(o)=\Phi(h\cdot o)=\Phi(o')=\lim_{t\to0}\Phi(g'(t)\cdot x')=\Phi(x').\]
So $\Phi$ is constant on fibres; the converse holds because orbits lie in fibres. The review record graded this clause as a
sketch; with the degeneration of \auxref{B.27}(c), whose reason is given in Appendix~\ref{app:auxiliary}, the steps above
complete it. Invariance alone does not suffice off the full-rank locus: on the invariant open set $\{\rk A=r\}$ the
projector $P_{\operatorname{row}A}$ is $\GL_r$-invariant, since $\operatorname{row}(gA)=\operatorname{row}A$, yet $(0,A)$
and $(0,A')$ with $\operatorname{row}A\ne\operatorname{row}A'$ lie in the fibre over $0$ and receive different values. This
$\Phi$ has no continuous extension to the whole factor space.

\emph{Task arithmetic.} $\sum_iw_is_iB_iA_i$ depends only on the products, so it is invariant, and
\[\sum_iw_is_iB_iA_i=\big[\,w_1s_1B_1\ \big|\ \cdots\ \big|\ w_Ns_NB_N\,\big]\begin{bmatrix}A_1\\ \vdots\\ A_N\end{bmatrix}\]
is a LoRA update of rank at most $\sum_ir_i$ with scale $1$, the concatenation \texttt{cat} of \thmref{Observation}{I.7}.

\emph{LoraHub.} Expanding, $\big(\sum_iw_iB_i\big)\big(\sum_jw_jA_j\big)=\sum_{i,j}w_iw_jB_iA_j$. Under the gauge
$(B_2,A_2)\mapsto(cB_2,A_2/c)$ the terms with $i,j\ne2$ and the term $w_2^2B_2A_2$ are unchanged, each $w_2w_lB_2A_l$ ($l\ne2$) gains the factor $c$,
and each $w_lw_2B_lA_2$ gains $c^{-1}$, which gives the displayed change. For $c\ne1$, nonzero weights and factors for which
the matrices $B_2A_l$, $B_lA_2$ ($l\ne2$) are linearly independent (true for generic factors), the change is nonzero. HF's
\texttt{linear} combination is $\big(\sum_i\beta_iB_i\big)\big(\sum_j\alpha_jA_j\big)$ with $\beta_i=\sqrt{\lvert w_is_i\rvert}$
and $\alpha_j=\operatorname{sign}(w_j)\sqrt{\lvert w_js_j\rvert}$ \citep{xmangrulkar2022peft}, a product of sums with the same
cross terms and the same failure. If all adapters share $A$, it equals $\big(\sum_j\alpha_j\big)\sum_i\beta_iB_iA$, which is a
fixed multiple of $\sum_iw_is_iB_iA$ for all $B_i$ only if the $\beta_i$ are proportional to the $w_is_i$, that is, only if all
$\lvert w_is_i\rvert$ with $w_is_i\ne0$ are equal (and the signs agree); for instance $w_1s_1=1$ and $w_2s_2=4$ give relative
weights $1:2$ instead of $1:4$.

\emph{Poly and MHR.} For task $t$ with routing weights $\alpha_{tj}$ over skills $j$ the update is
$W_t=\big(\sum_j\alpha_{tj}B_j\big)\big(\sum_j\alpha_{tj}A_j\big)$. A common $g$ gives
$\big(\sum_j\alpha_{tj}B_jg^{-1}\big)\big(\sum_j\alpha_{tj}gA_j\big)=W_t$, so the diagonal $\GL_r$ preserves the map. Per-skill
gauges $(g_j)$ that preserve it for all routings in an open set are diagonal on any pair of mixed skills with complementary
column spaces: $W_t=\sum_{i,j}\alpha_{ti}\alpha_{tj}B_iA_j$ is a polynomial in $\alpha$, and comparing the coefficients of
$\alpha_{ti}\alpha_{tj}$, $i\ne j$, gives
\[B_i(g_i^{-1}g_j-I)A_j+B_j(g_j^{-1}g_i-I)A_i=0 .\]
If $\operatorname{col}B_i\cap\operatorname{col}B_j=0$, the two terms vanish separately, and for injective $B_i$ and surjective
$A_j$ this forces $g_i=g_j$. That for one fixed generic routing with fewer skills than tasks the continuous fibre is exactly
the diagonal orbit is the numerical check reported below. With at least as many skills as tasks the fibre is larger. If there
are more skills than tasks, a vector $v\ne0$ with $\sum_jv_j\alpha_{tj}=0$ for every $t$ exists, and $B_j\mapsto B_j+v_jX$
leaves every $\sum_j\alpha_{tj}B_j$ unchanged for every $X$. If there are equally many and the routing matrix is invertible,
the task-wise sums $(C_t,D_t)=\big(\sum_j\alpha_{tj}B_j,\sum_j\alpha_{tj}A_j\big)$ are free coordinates, and each task has its
own $\GL_r$ gauge. Counting the routing as a parameter, relabelling the skills together with the columns of $\alpha$ preserves
the map, and so does $(\alpha,B)\mapsto(\alpha M,M^{-1}B)$ for invertible $M$ with $M\mathbf 1=\mathbf 1$, where $B$ stacks the
skills (and likewise for $A$), since $(\alpha M)(M^{-1}B)=\alpha B$ and $M\mathbf 1=\mathbf 1$ keeps the rows of $\alpha$ summing
to $1$. Poly and MHR train skills and routing jointly through this map, so it is the method's own $\rho$; the question of
(a) concerns combining modules that were factorised independently.

(b) \emph{Gauge fixing.} If a section $\sigma$ of $\pi$ chooses the factors from the products, then $\Phi\circ\sigma$ is a
function of the products. \emph{Exactness.} Task arithmetic $T(w)=\sum_iw_iB_iA_i$ (scales absorbed) is linear in $w$.
Consider rank-preserving products of sums in which each adapter's factor carries its own weight $w_i$, in one sum or in
both. If both sums are weighted, as in LoraHub, the rule $P(w)$ is homogeneous of degree two in $w$, and $P(w)=T(w)$ for all
$w$ would give, replacing $w$ by $\lambda w$, $\lambda^2P(w)=\lambda T(w)$ for all $\lambda$, hence $T\equiv0$. So an exact
rule of this kind has degree one in $w$: one of its factors does not depend on $w$, say $\big(\sum_iw_iB_i\big)X$. Then exactness for all $w$ means $B_iX=B_iA_i$ for every $i$, and for injective $B_i$ this gives
$A_i=X$: the $A$-factor is shared. The symmetric case $X\big(\sum_jw_jA_j\big)$ with surjective $A_j$ forces a shared $B$.
\texttt{cat} realises $T(w)$ exactly, as shown in (a), with inner dimension $\sum_ir_i$. With a frozen shared $A$,
$\big(\sum_iw_iB_i\big)A=\sum_iw_iB_iA$, and LoraHub's rule gives $\big(\sum_iw_iB_i\big)\big(\sum_jw_j\big)A$.

\emph{SVD canonical forms.} If $B_iA_i=U_i\Sigma_iV_i\T$ and the canonical factors are $B_i=U_i\Sigma_i^{1/2}$,
$A_i=\Sigma_i^{1/2}V_i\T$ (or any convention built from $U_i,\Sigma_i,V_i$), then changing the signs of the $j$-th column of
$U_i$ and of $V_i$ gives another SVD and another canonical pair, multiplying column $j$ of $B_i$ and row $j$ of $A_i$ by $-1$.
The product $B_iA_i$ is unchanged, while a cross term $B_iA_l$ ($l\ne i$) changes by $-2\,b_{ij}a_{lj}\T$, where $b_{ij}$ is
column $j$ of $B_i$ and $a_{lj}\T$ row $j$ of $A_l$, which is nonzero generically. When $\sigma_j=\sigma_{j+1}$, rotating the
corresponding columns of $U_i$ and $V_i$ by the same $O\in\Ort(2)$ also preserves the product and changes the cross terms.
\emph{KnOTS.} Write the horizontal concatenation $[\Delta_1\mid\cdots\mid\Delta_N]=U\Sigma[V_1\T\mid\cdots\mid V_N\T]$ by an
SVD \citep{stoica2024model}; then $\Delta_i=(U\Sigma)V_i\T$ share the left factor $U\Sigma$, and
$(U\Sigma)\big(\sum_iw_iV_i\T\big)=\sum_iw_i\Delta_i$ exactly. \emph{FedAvg-LoRA.} Expanding,
\[\tfrac12(B_1+B_2)\cdot\tfrac12(A_1+A_2)-\tfrac12(B_1A_1+B_2A_2)=-\tfrac14(B_1-B_2)(A_1-A_2),\]
which is nonzero whenever $(B_1-B_2)(A_1-A_2)\ne0$, whatever gauge the broadcast pair is in.

(c) Write a merge rule as $\Psi(\tau_1,\dots,\tau_K)$ for task vectors $\tau_k\in\R^N$ (or the matrices with $N$ entries,
flattened). It is $g$-equivariant if $\Psi(g\tau_1,\dots,g\tau_K)=g\Psi(\tau_1,\dots,\tau_K)$, read for random rules either
for each realisation of the randomness (per mask) or as equality of laws (in distribution).

\emph{Task arithmetic} is linear, hence $\GL(\Theta)$-equivariant.

\emph{DARE.} $\mathrm{DARE}(\tau)=m\odot\tau/(1-p)$, with a mask $m$ of independent Bernoulli$(1-p)$ entries, $0<p<1$,
drawn independently for each task, followed by a weighted sum. For invertible diagonal $D$, $D(m\odot\tau)=m\odot(D\tau)$, so
$\mathrm{DARE}(D\tau)=D\,\mathrm{DARE}(\tau)$ for each mask, and the weighted sum preserves this. For a permutation matrix $P$,
$P(m\odot\tau)=(Pm)\odot(P\tau)$ and $Pm$ has the law of $m$, so $\mathrm{DARE}(P\tau)$ has the law of $P\,\mathrm{DARE}(\tau)$,
jointly over the tasks. Together, $\mathrm{Mon}_N=S_N\ltimes T_N$ acts equivariantly in distribution. Conversely, suppose
$\mathrm{DARE}(g\tau)$ and $g\,\mathrm{DARE}(\tau)$ have the same law for every $\tau$, and take $\tau=e_k$. The law of
$g\,\mathrm{DARE}(e_k)$ is supported on the two points $0$ and $ge_k/(1-p)$, while $\mathrm{DARE}(ge_k)$ takes the $2^q$
distinct values $m\odot ge_k/(1-p)$, each with positive probability, where $q$ is the number of nonzero entries of $ge_k$. So
$q=1$: every column of $g$ has a single nonzero entry, and $g\in\mathrm{Mon}_N$.

\emph{TIES with global trimming} \citep{yadav2023ties}. Each $\tau_k$ is trimmed to its $k_0$ entries of largest magnitude,
giving $\hat\tau_k$; the sign $\gamma_i=\operatorname{sign}\big(\sum_k\hat\tau_{k,i}\big)$ is elected by total mass; and the
output is $\lambda$ times the disjoint mean, $\frac1{\lvert A_i\rvert}\sum_{k\in A_i}\hat\tau_{k,i}$ over
$A_i=\{k:\operatorname{sign}\hat\tau_{k,i}=\gamma_i\}$. Let $g=cPS$ with $c\ne0$, $P$ a permutation matrix and $S=\diag(s_1,\dots,s_N)$ a
diagonal sign matrix, and let $j$ be the coordinate that $P$ moves to $i$. The magnitudes of $g\tau_k$ are $\lvert c\rvert$ times
the permuted magnitudes of $\tau_k$, so, off ties at the threshold, trimming keeps the corresponding entries and
$\widehat{g\tau_k}=g\hat\tau_k$. The elected sign at $i$ becomes $\operatorname{sign}(c)s_j\gamma_j$, the set $A_i$ for the
transformed input is $A_j$, and the mean at $i$ is $cs_j$ times the mean at $j$. So the output is $g$ times the output, and
TIES is $\R^\times\cdot B_N$-equivariant off ties. A non-scalar positive diagonal fails: with $N=2$, one task, $k_0=1$,
$\tau=(2,1)$ and $D=\diag(1,3)$, TIES returns $\lambda(2,0)$ on $\tau$ and $\lambda(0,3)\ne D\lambda(2,0)$ on $D\tau$.

\emph{Frequency election.} Here $\gamma_i=\operatorname{sign}\big(\sum_k\operatorname{sign}\hat\tau_{k,i}\big)$, with $0$
replaced by $+1$. Positive homotheties and permutations leave the counts unchanged or permute them, so the argument above
gives $\R_{>0}\cdot S_N$-equivariance. Sign changes fail on ties, which have positive probability because the counts are
integers. With two tasks and $\hat\tau_{1,i}=1$, $\hat\tau_{2,i}=-1$, the tie elects $+1$ and the output at $i$ is $\lambda$;
after flipping coordinate $i$ the entries are $-1,+1$, the tie again elects $+1$, and the output at $i$ is $\lambda$, not
$-\lambda$. The same example defeats $-I$.

\emph{Per-tensor trimming.} If coordinates are grouped into tensors $b$ of sizes $N_b$ and each tensor is trimmed to its own
top entries, an independent homothety $c_b$ and signed permutation inside each tensor preserve the trimmed sets, and
permuting tensors of equal size (and equal $k_0$) preserves them too; the election and the mean are coordinatewise, so the
argument above applies blockwise.

\emph{DARE followed by TIES.} For $g\in\R^\times\cdot B_N\subseteq\mathrm{Mon}_N$, the masked inputs
$(\mathrm{DARE}(g\tau_k))_k$ have the law of $(g\,\mathrm{DARE}(\tau_k))_k$, and TIES commutes with $g$ off magnitude ties
among nonzero entries (ties among entries zeroed by the mask do not matter, since a kept zero is zero), so the composite is
equivariant in distribution. HF's and MergeKit's \texttt{dare\_ties} \citep{xmangrulkar2022peft,xgoddard2024mergekit} omit the
trim: for invertible diagonal $D$ and fixed masks, $\mathrm{DARE}(D\tau_k)=D\,\mathrm{DARE}(\tau_k)$, the elected sign at $i$
is multiplied by $\operatorname{sign}d_i$, the sets $A_i$ are unchanged and the mean is multiplied by $d_i$; so it is exactly
$T_N$-equivariant per mask, and with permutations in distribution its group contains $\mathrm{Mon}_N$.

\emph{No rotations.} Let $R$ rotate the plane of $e_1,e_2$ by $30^\circ$, so $Re_1=(\cos30^\circ,\sin30^\circ,0,\dots)$. For DARE,
$R\,\mathrm{DARE}(e_1)$ takes the values $0$ and $Re_1/(1-p)$, each with two nonzero entries or none, while
$\mathrm{DARE}(Re_1)$ has exactly one nonzero entry with probability $2p(1-p)>0$; the laws differ. For TIES with one task and
$k_0=1$, $\mathrm{TIES}(e_1)=\lambda e_1$, so $R\,\mathrm{TIES}(e_1)$ has two nonzero entries, while $\mathrm{TIES}(Re_1)=
\lambda\cos30^\circ e_1$ has one. With one task, \texttt{dare\_ties} reduces to DARE, and DARE followed by trimmed TIES
returns, with positive probability, the trimmed vector $\lambda\cos30^\circ e_1/(1-p)$ on $Re_1$, which no value of
$R\,\mathrm{TIES}(\mathrm{DARE}(e_1))\in\{0,\lambda Re_1/(1-p)\}$ matches.

\emph{Factor-level rules.} HF's non-SVD variants apply the sparsifier to each $B_i$ and each $A_j$ with independent masks and
then form $\big(\sum_iv_i\,\mathrm{DARE}(B_i)\big)\big(\sum_jv_j\,\mathrm{DARE}(A_j)\big)$ (and the analogues for TIES). The
$B$-masks and the $A$-masks are independent and DARE is unbiased, $\mathbb E\,\mathrm{DARE}(X)=X$, so the expected output is the
LoraHub-type rule $\big(\sum_iv_iB_i\big)\big(\sum_jv_jA_j\big)$. A gauge $(cB_2,A_2/c)$ on one adapter changes this expectation
by the formula of (a), which is nonzero for generic factors, so the output's law changes: even DARE is not invariant under
independent per-adapter gauges, not even in distribution. On the other hand, let $g\in T_r$ act on all adapters at once,
$B_i\mapsto B_ig^{-1}$ and $A_j\mapsto gA_j$. Right multiplication by a diagonal matrix scales columns, which commutes with an
entrywise mask, so $M\odot(B_ig^{-1})=(M\odot B_i)g^{-1}$ and $M'\odot(gA_j)=g(M'\odot A_j)$, and for every choice of masks
\[\Big(\sum_iv_i\,M_i\odot B_ig^{-1}\Big)\Big(\sum_jv_j\,M_j'\odot gA_j\Big)=\Big(\sum_iv_i\,M_i\odot B_i\Big)g^{-1}g\Big(\sum_jv_j\,
M_j'\odot A_j\Big).\]
For a permutation $g=P$, $M\odot(B_iP\T)=\big((MP)\odot B_i\big)P\T$, and $MP$ has the law of $M$; with all masks permuted by the
same $P$ their joint law is unchanged, so the output's law is invariant. A single masked adapter is the case of one
summand. This gives $T_r$-invariance per mask and $\mathrm{Mon}_r$-invariance in distribution.
\end{proof}

\begin{remark*}{numerical checks}
The checks reported after \thmref{Proposition}{V.4}, namely DARE's per-mask equivariance under a random diagonal to
$2\times10^{-15}$, TIES's homogeneity and its failure by $5.6$ under a positive diagonal, the LoraHub formula with three
modules, the order-one change of factor-wise DARE under a one-adapter gauge against $10^{-15}$ under a shared one, and the Poly
fibre dimensions $32$ (two tasks) and $4=\dim\GL_2$ (four tasks) for three skills, $r=2$ and $6\times6$ weights, instantiate the
statements proved above; the last one is the numerical part of the Poly bullet (Appendix~\ref{app:repro}).
\end{remark*}

% ====================================================================================================
\subsection{Definition VI.1}\label{proof:VI.1}

\begin{restate}{Definition}{VI.1}{extensions; mergeability}
\begin{itemize}
  \item An \textbf{extension} is a function-preserving pointed embedding
  \[\iota:(\Theta,f,\theta_0)\to(\Theta\times\Theta',\tilde f,(\theta_0,\theta_0')),\qquad\iota(\theta)=(\theta,\theta_0'),\]
  with $\tilde f(\theta,\theta_0',-)=f(\theta,-)$. The point $\theta_0'$ is the \textbf{neutral point} (``identity at
  initialisation'').
  \item An \textbf{unpointed extension} has no neutral point.
  \item A method on an extension is \textbf{mergeable} if there is $\mu:Q\to\Theta$ with
  $\Phi_{\tilde f}\circ\rho=\Phi_f\circ\mu$: a lift along the realisation.
  \item It is \textbf{box-locally mergeable} if the modified box function stays in the box's function class: linear maps
  for a linear box, affine maps when the box has a bias slot. \thmref{Theorem}{VI.3} gives sufficient conditions, with
  explicit side conditions (position-uniformity, no weight sharing, bias slots), under which this yields a merge into the
  base network.
\end{itemize}
\end{restate}

A definition needs no proof. We check the facts that the definition and the square after it use. Here
$\Phi_f:\Theta\to C^\infty(X,Y)$, $\theta\mapsto f(\theta,-)$, is the realisation, and a box's function class consists of the
maps it can compute with some weight, applied identically at every position.

\begin{proof}[Verification of the claims in Definition VI.1]
\emph{Extensions.} The map $\iota(\theta)=(\theta,\theta_0')$ is smooth, injective and immersive, and a homeomorphism onto the
closed set $\Theta\times\{\theta_0'\}$, so it is an embedding. It is pointed, $\iota(\theta_0)=(\theta_0,\theta_0')$, and
function-preserving, $\Phi_{\tilde f}\circ\iota=\Phi_f$, which is the displayed condition.

\emph{Weight-space methods merge.} A method $(Q,q_0,\rho)$ on $\Theta$ itself, read on the trivial extension
($\Theta'$ a point), satisfies the merge condition with $\mu=\rho$, the arrow to full fine-tuning of
\thmref{Observation}{I.4}(a). More generally, a method on an extension whose $\rho$ factors as $\iota\circ\mu$ merges by
$\mu$.

\emph{Merges are lifts, and need not be unique.} The realisation $\Phi_f$ need not be injective. Every function-preserving
reparametrisation $\phi$ of the base gives another lift $\phi\circ\mu$; permutations of hidden units are examples
(\thmref{Lemma}{VI.2}).
Let a pointed method merge, with $\rho(q_0)$ the neutral point $(\theta_0,\theta_0')$. Then
\[\Phi_f(\mu(q_0))=\Phi_{\tilde f}(\theta_0,\theta_0')=\Phi_f(\theta_0),\]
so $\mu(q_0)$ realises the pretrained function, without having to equal $\theta_0$.

\emph{Box-local mergeability.} If the edit of a method modifies a single linear (affine) box of the base network, and for
every $q$ the modified box function is linear (affine), position-independent and computed by an unshared weight $W(q)$ of
that box, then replacing the box's weight by $W(q)$ and keeping the others gives a $\mu$ with $\Phi_f\circ\mu=\Phi_{\tilde
f}\circ\rho$. The side conditions are exactly what makes this replacement legitimate, and \thmref{Theorem}{VI.3}(a)
spells them out.
\end{proof}

% ====================================================================================================
\subsection{Lemma VI.2}\label{proof:VI.2}

\begin{restate}{Lemma}{VI.2}{the activation fixes the privileged basis; classical, \citet{xsussmann1992uniqueness},
\citet{xchen1993geometry}, \citet{xgodfrey2022symmetries}}
Let $\sigma:\R\to\R$ be an elementwise activation that is not affine, and let $L\in\GL_d$ commute with it as maps on $\R^d$:
$L\,\sigma^{\times d}=\sigma^{\times d}L$.

(a) If $\sigma\in C^2$, then $L=PD$ is monomial.

(b) If $\sigma$ is ReLU or leaky ReLU with slope $\alpha\ne1$, then $L=PD$ with $D$ a positive diagonal.

(c) If moreover $\sigma$ is real-analytic at $0$ with $\sigma(t)=\sum a_kt^k$, each diagonal entry $c$ of $D$ satisfies
$ca_k=c^ka_k$ for all $k$. So $c=1$ for GELU and SiLU, since $\sigma''(0)\neq0$, and $c=\pm1$ for tanh.
\end{restate}

\begin{proof}
Write $\ell_i$ for row $i$ of $L$; it is nonzero because $L$ is invertible. Row $i$ of the identity
$L\,\sigma^{\times d}(x)=\sigma^{\times d}(Lx)$ reads
\begin{equation}\label{eq:C3-commute-row}
\sum_jL_{ij}\,\sigma(x_j)=\sigma(\ell_i\cdot x)\qquad\text{for all }x\in\R^d .
\end{equation}
Once each row has a single nonzero entry, the positions of these entries are distinct, since otherwise a column of $L$
would vanish; so $L=PD$ with $P$ a permutation matrix and $D$ an invertible diagonal matrix.

(a) For $j\ne k$, the mixed partial $\partial_j\partial_k$ of the left side of \eqref{eq:C3-commute-row} vanishes, since each
term depends on one coordinate, and the mixed partial of the right side is $L_{ij}L_{ik}\,\sigma''(\ell_i\cdot x)$. Since
$\sigma$ is $C^2$ and not affine, $\sigma''(t_0)\ne0$ for some $t_0$, and since $\ell_i\ne0$ there is $x$ with $\ell_i\cdot
x=t_0$. Hence $L_{ij}L_{ik}=0$ for all $j\ne k$: row $i$ has a single nonzero entry.

(b) Let $\sigma(t)=t$ for $t\ge0$ and $\sigma(t)=\alpha t$ for $t<0$, with $\alpha\ne1$ ($\alpha=0$ is ReLU). Then $\sigma$ is
smooth off $0$, and not differentiable at $0$, where its one-sided derivatives $1$ and $\alpha$ differ. The left side of
\eqref{eq:C3-commute-row} is a sum of functions of separate coordinates, so it is differentiable at $x$ iff each term with
$L_{ij}\ne0$ is differentiable in $x_j$; its non-differentiability set is $\bigcup_{j:L_{ij}\ne0}\{x_j=0\}$. The right side is
smooth where $\ell_i\cdot x\ne0$, and at a point with $\ell_i\cdot x=0$ its one-sided directional derivatives along $\pm\ell_i$
are $\lVert\ell_i\rVert^2$ and $-\alpha\lVert\ell_i\rVert^2$, which are not negatives of each other, so it is not
differentiable there; its non-differentiability set is the hyperplane $\{\ell_i\cdot x=0\}$. The two sets are equal. By
\auxref{B.36}(b) the hyperplane is one of the coordinate hyperplanes $\{x_{j_0}=0\}$, and each $\{x_j=0\}$ with $L_{ij}\ne0$
is contained in it and so equal to it; hence $L_{ij}\ne0$ only for $j=j_0$. Then \eqref{eq:C3-commute-row} reads
$L_{ij_0}\sigma(t)=\sigma(L_{ij_0}t)$ for all $t$. At $t=1$ this is $L_{ij_0}=\sigma(L_{ij_0})$; if $L_{ij_0}<0$ the right side
is $\alpha L_{ij_0}$, which would force $\alpha=1$. So $L_{ij_0}>0$, and $D$ is positive.

(c) A permutation matrix commutes with every elementwise map (rule R2 of \thmref{Theorem}{VI.3}), so
$D\sigma^{\times d}=P^{-1}L\sigma^{\times d}=P^{-1}\sigma^{\times d}L=\sigma^{\times d}P^{-1}L=\sigma^{\times d}D$, that is,
$c\,\sigma(t)=\sigma(ct)$ for each diagonal entry $c\ne0$ and all $t$. Near $0$ both sides are power series,
$\sum_kca_kt^k=\sum_ka_kc^kt^k$, so $ca_k=c^ka_k$ for every $k$ (\auxref{B.10}(c)). For GELU, $\sigma(t)=t\,\Phi_{\mathcal
N}(t)$ with the standard normal distribution function $\Phi_{\mathcal N}(t)=\tfrac12+t/\sqrt{2\pi}+O(t^3)$, so
$a_2=1/\sqrt{2\pi}\ne0$; for SiLU, $\sigma(t)=t/(1+e^{-t})=\tfrac t2+\tfrac{t^2}4+O(t^4)$, so $a_2=\tfrac14\ne0$. In both cases
$c=c^2$, and $c\ne0$ gives $c=1$. For tanh, $a_3=-\tfrac13\ne0$, so $c=c^3$ and $c=\pm1$; both occur, tanh being odd.
\end{proof}

% ====================================================================================================
\subsection{Theorem VI.3}\label{proof:VI.3}

\begin{restate}{Theorem}{VI.3}{the merge calculus; known in part}
Work with activation-side diagrams in a symmetric monoidal category of continuous, piecewise-smooth maps with copy and
discard nodes. (We call copy and the Hadamard product ``spiders'' by analogy only; in this cartesian setting they do not
satisfy the Frobenius law.) A \emph{trainable box} is a node whose weight ranges over a stated class: all linear maps, or all
affine maps when the box has a bias slot. R2 and R4--R7 are identities of maps; R1 is a closure property and R3 a
characterisation.

\smallskip
{\upshape\footnotesize\renewcommand{\arraystretch}{1.15}\noindent
\begin{tabularx}{\linewidth}{@{}>{\raggedright\arraybackslash}p{0.17\linewidth}>{\raggedright\arraybackslash}p{0.33\linewidth}>{\raggedright\arraybackslash}X@{}}
\toprule\headrow
\textbf{Rule} & \textbf{Statement} & \textbf{When it holds}\\
\midrule
R1 (fusion) & composites and parallel sums of linear (affine) boxes are linear (affine) & the operators involved are
position-independent\\
R2 (permutation slide) & $P\sigma^{\times d}=\sigma^{\times d}P$ & every elementwise $\sigma$\\
R3 (positive-diagonal slide) & $\diag(\ell)\,\sigma^{\times d}=\sigma^{\times d}\diag(\ell)$ for all $\ell\ge0$ & iff
$\sigma$ is positively homogeneous (ReLU, leaky ReLU; not GELU, SiLU, tanh); such $\sigma$ are piecewise linear\\
R4 (spider slide) & $\diag(\ell)(a\odot b)=(\diag(\ell)a)\odot b$ & any $\ell$\\
R5 (cup slide) & $\langle Mq,k\rangle=\langle q,M^\top k\rangle$ & always. Under RoPE,
$\langle R_mMq,R_nk\rangle=\langle q,M^\top R_{n-m}k\rangle$ holds for every $M$, but the slide onto the key side,
$\langle R_mMq,R_nk\rangle=\langle R_mq,R_nM^\top k\rangle$ for all $m,n$, holds iff $M^\top$ commutes with every
$R_{n-m}$ that occurs; for diagonal $M=\diag(\ell)$ and distinct frequencies, iff $\ell$ is constant on each rotary pair
(in either pairing convention)\\
R6 (norm absorption) & $W(\gamma\odot n(x)+\beta)=(W\diag\gamma)\,n(x)+W\beta$ & always; the shift $W\beta$ needs a bias
slot, and an RMSNorm gain has no $\beta$\\
R7 (copy slide) & applying $D$ and then copying equals copying and then applying $D$ to every copy & always\\
\bottomrule
\end{tabularx}}
\smallskip

Consequences:

(a) \textbf{Affine edits.} Let a method insert a position-uniform affine edit $E_q$ on a wire $w$ of the base network $f$,
for every $q\in Q$. Suppose every reader of $w$ is a trainable, unshared box (no weight tying) whose class is closed under
precomposition with $E_q$; a shift needs a bias slot. Then the method merges, by $W_c\mapsto W_cE_q$ for all readers $c$ at
once. Dually, it merges into an untied trainable box that writes $w$ whenever $E_q$ composed with that box stays in its
class. Absorption into a box that the extension itself inserts does not count. Edits that violate these hypotheses fall
outside (a): on the reader side, R1--R7 do not merge position-selective edits or edits on a wire that is also read by a skip
connection, a norm or the network output. Whole-network non-mergeability is open. Position-selective edits (ReFT/LoReFT, and
position-gated LoRA variants such as Activated LoRA) are not box-locally mergeable. Residual-stream edits need case
analysis: an edit merges into the box that produces the wire when that box can absorb it (an untied embedding table; for a
fixed-position edit at layer $0$, a learned absolute positional table), a shift folds into a writer's bias when one exists,
a diagonal edit after a post-LN folds into that LN's affine parameters, and an orthogonal edit in an RMSNorm model merges
globally, SliceGPT/QuaRot-style. No merge is known for LoReFT's non-orthogonal low-rank edit.

(b) \textbf{Diagonal methods.} A diagonal method merges into a box reached by R1--R7 if all of the following hold: the rules
move \emph{every} copy of $\diag(\ell)$ into trainable boxes for all $\ell$ in the method's parameter domain (with $\ell\ge0$
when R3 is used); no position-dependent operator lies on the path (for RoPE, $\ell$ is constant on each rotary pair); the
target box is unshared, or $\ell$ is constant across the sharing group; and every shift has a bias slot. Conversely, the
RoPE condition is necessary for sliding a query scaling onto $W_K$ when $W_Q$ and $W_K$ have full rank and at least two
occurring relative offsets have rotation angles that are not multiples of $\pi$. Without such genericity hypotheses the
conditions are not necessary: a dead rotary pair, or zero entries of $\ell$ under ReLU, can still merge.

(c) \textbf{Non-affine edits.} For a linear or affine box, an edit whose modified box function is non-affine in the box
input, for some parameter value, is not box-locally mergeable. For linear and affine boxes this is the definition; for
non-linear boxes such as attention, see \thmref{Theorem}{VI.5}(f).

(d) \textbf{Type relativity.} A generic high-precision update is off the quantisation grid, so exact merging of a QLoRA or
LoftQ update leaves the quantised type, a type-preserving merge is lossy, and QA-LoRA's group-constant adapters are the
exception (\thmref{Proposition}{V.2}(d)).
\end{restate}

\begin{proof}
A diagram denotes a map on activations, and two diagrams related by an identity of maps compute the same network function;
a rewrite that ends with every edit inside trainable boxes of the base network therefore exhibits a merge $\mu$, the new
weights of those boxes. ``Position-independent'' means that the same map acts at every token position.

\emph{R1.} Composites and sums of linear (affine) maps are linear (affine). If each acts identically at every position, so
does the result, which is then computed by one weight (and bias) at all positions.

\emph{R2.} For a permutation $\pi$ with matrix $P$, $(P\sigma^{\times d}(x))_i=\sigma(x_{\pi^{-1}(i)})=\sigma((Px)_i)$.

\emph{R3.} Coordinatewise the identity says $\ell_i\sigma(t)=\sigma(\ell_it)$ for all $\ell_i\ge0$ and $t$; with $d=1$ this is
positive homogeneity of degree one, and conversely positive homogeneity gives the identity in every coordinate. A positively
homogeneous $\sigma$ satisfies $\sigma(t)=t\,\sigma(1)$ for $t\ge0$ and $\sigma(t)=\lvert t\rvert\,\sigma(-1)$ for $t<0$, so it
is linear on each half-line. ReLU and leaky ReLU are of this form. GELU, SiLU and tanh are differentiable at $0$ and not
linear, so by \auxref{B.36}(a) they are not positively homogeneous.

\emph{R4} holds entrywise, $\ell_ia_ib_i=(\ell_ia_i)b_i$.

\emph{R5.} $\langle Mq,k\rangle=q\T M\T k=\langle q,M\T k\rangle$. Under RoPE, $R_m$ is block-diagonal with the planar rotations
$\rho(m\omega_j)$ on the rotary pairs, in either pairing convention, so $R_m\T=R_{-m}$ and $R_m\T R_n=R_{n-m}$. Hence
$\langle R_mMq,R_nk\rangle=q\T M\T R_{n-m}k=\langle q,M\T R_{n-m}k\rangle$ for every $M$, while
$\langle R_mq,R_nM\T k\rangle=\langle q,R_{n-m}M\T k\rangle$. These agree for all $q,k$ iff $M\T R_{n-m}=R_{n-m}M\T$, and for all
occurring $(m,n)$ iff $M\T$ commutes with every occurring $R_\Delta$. For $M=\diag(\ell)$ the commutation holds pair by pair.
On a pair with entries $(\ell_a,\ell_b)$ and rotation angle $\varphi$,
\[\begin{pmatrix}\ell_a&0\\0&\ell_b\end{pmatrix}\begin{pmatrix}\cos\varphi&-\sin\varphi\\ \sin\varphi&\cos\varphi\end{pmatrix}
-\begin{pmatrix}\cos\varphi&-\sin\varphi\\ \sin\varphi&\cos\varphi\end{pmatrix}\begin{pmatrix}\ell_a&0\\0&\ell_b\end{pmatrix}
=(\ell_b-\ell_a)\sin\varphi\begin{pmatrix}0&1\\1&0\end{pmatrix},\]
which vanishes iff $\ell_a=\ell_b$ or $\varphi\in\pi\mathbb Z$. So the slide holds iff $\ell$ is constant on every pair whose
angle $\Delta\omega_j$ is not a multiple of $\pi$ for some occurring $\Delta$. For RoPE, $0<\omega_j\le1$, so every pair
qualifies as soon as the offset $\Delta=1$ occurs; distinct frequencies are not needed for diagonal $M$.

\emph{R6.} By linearity, $W(\gamma\odot n(x)+\beta)=W\diag(\gamma)\,n(x)+W\beta$. The first term is a linear box applied to the
normalised wire, and the constant $W\beta$ is absorbed by a bias slot, $b\mapsto b+W\beta$; an RMSNorm gain has $\beta=0$.

\emph{R7.} With $\Delta(x)=(x,x)$ the copy map, $\Delta\circ D=(D\times D)\circ\Delta$ for every map $D$: the naturality of the
cartesian diagonal (\bgref{B.4}).

(a) Let $E_q(x)=A_qx+c_q$ act at every position on $w$. By R7 applying $E_q$ before the copy node that feeds the readers
equals applying it on each reader's input. For a reader $c$ with weight $W_c$ and bias $b_c$,
$W_c(A_qx+c_q)+b_c=(W_cA_q)x+(W_cc_q+b_c)$, which lies in the class of $c$ by the closure hypothesis (for a linear class this
requires $c_q=0$, and for a shift a bias slot). The readers are unshared, so their weights can be changed independently; and
$E_q$ is the same at every position, so one weight serves all positions (R1). Setting $\mu(q)$ to the base weights with
$W_c\mapsto W_cA_q$ and $b_c\mapsto b_c+W_cc_q$ for every reader gives $\Phi_f\circ\mu=\Phi_{\tilde f}\circ\rho$. Dually, if
$w$ is the output of an untied box $p(x)=W_px+b_p$ and $E_q\circ p$ lies in its class, replace $p$ by $E_q\circ p$, that is,
$W_p\mapsto A_qW_p$ and $b_p\mapsto A_qb_p+c_q$ (R1). A merge must land in the base network $\Theta$, so absorption into a box
that only the extension contains is not a merge. On the reader side the rules have nothing to absorb an edit read by a skip
connection, a norm or the network output, since those readers carry no trainable weight, and R1 does not fuse a
position-selective edit into a weight that acts at every position. A position-selective edit modifies the box function at
some positions only; for parameter values at which it is not the identity, the modified computation is not one map applied
identically at every position, so it is not in the box's class, and the edit is not box-locally mergeable.

\emph{Residual-stream cases.} (i) If $w$ is the output of an untied embedding table $\mathcal E$, the dual case applies:
replacing each row $e(v)$ by $A_qe(v)+c_q$ absorbs a position-uniform affine edit. For an additive edit $h\mapsto h+c$ at one
fixed position $p_0$ of layer $0$ in a model with a learned absolute positional table, replacing row $p_0$ of that table by
itself plus $c$ absorbs it. (ii) A shift $h\mapsto h+t$ of the stream after a block whose last box has a bias $b$ is absorbed
by $b\mapsto b+t$, since the block adds its output to the stream. (iii) In a post-LN model the normalised wire is the stream,
$h=\gamma\odot n(x)+\beta$, and a diagonal edit gives $s\odot h+t=(s\odot\gamma)\odot n(x)+(s\odot\beta+t)$, absorbed by
$\gamma\mapsto s\odot\gamma$, $\beta\mapsto s\odot\beta+t$. (iv) In an untied pre-RMSNorm model, let block $j$ map the stream by
$h\mapsto h+F_j(h)$ with $F_j(h)=W^{(j)}_{\rm out}\,G_j\big(W^{(j)}_{\rm in}\diag(\gamma_j)\,n(h)\big)$, where $n(h)=h/\mathrm{rms}(h)$,
$G_j$ is the internal computation (attention with its positional encoding applied after the input projections, or the gated
MLP), and $W^{(j)}_{\rm in}$ collects the input projections. Let $Q$ be orthogonal and the edit be $h\mapsto Qh$ after block
$l$. Replace the embedding rows $e(v)$ by $Qe(v)$, and for $j\le l$ replace $W^{(j)}_{\rm in}$ by $W^{(j)}_{\rm in}\diag(\gamma_j)Q\T$,
$\gamma_j$ by $\mathbf 1$ and $W^{(j)}_{\rm out}$ by $QW^{(j)}_{\rm out}$ (and any bias $b$ of an output projection by $Qb$). Since
$\mathrm{rms}(Qh)=\mathrm{rms}(h)$, $n(Qh)=Qn(h)$, and by induction the new streams are $Qh_j$ for $j\le l$:
$Qh_{j-1}+QW^{(j)}_{\rm out}G_j\big(W^{(j)}_{\rm in}\diag(\gamma_j)Q\T Qn(h_{j-1})\big)=Q\big(h_{j-1}+F_j(h_{j-1})\big)$. After block $l$
the new network carries $Qh_l$, as the edited one does, and the later blocks are unchanged, so the outputs agree. This is the
construction of \citet{xashkboos2024slicegpt,xashkboos2024quarot}.

(b) The rules used along the path are identities of maps under their side conditions: R3 needs $\ell\ge0$, the slide of R5
under RoPE needs $\ell$ constant on the rotary pairs, and R6 needs a bias slot for the shift. Applying them moves every copy of
$\diag(\ell)$ next to a trainable box, and R1 fuses it; this is (a) applied at the box reached. An unshared target can absorb
its copy independently, and a target shared across a group absorbs a single matrix, which is possible for every member exactly
when $\ell$ is constant across the group. For the converse under RoPE, suppose the query scaling $D=\diag(\ell)$ merges into
some $W_K'$ with $W_Q$ unchanged. Equality of all logits for all inputs gives
$x\T W_Q\T DR_\Delta W_Kx'=x\T W_Q\T R_\Delta W_K'x'$ for every occurring offset $\Delta=n-m$, so
$W_Q\T(R_\Delta W_K'-DR_\Delta W_K)=0$. If $W_Q$ has full row rank, $W_Q\T$ is injective and $W_K'=R_{-\Delta}DR_\Delta W_K$ for
every occurring $\Delta$; if $W_K$ has full row rank too, it has a right inverse, and $R_{-\Delta}DR_\Delta$ is one matrix for
all occurring $\Delta$. On a rotary pair $j$, let $\Delta$ and $\Delta'$ be occurring offsets with $(\Delta-\Delta')\omega_j\notin\pi
\mathbb Z$; then $D$ commutes with $R_{\Delta-\Delta'}$ on that pair, and by the computation in R5, $\ell$ is constant on it. In
causal attention $\Delta'=0$ occurs (a query attends to its own position), so it suffices that one occurring offset has an angle
on pair $j$ that is not a multiple of $\pi$, which the hypothesis provides pair by pair. The genericity hypotheses are needed.
If $W_Q$ has a dead pair (both rows zero) and $\ell$ is constant on every other pair, then $W_Q\T DR_\Delta=W_Q\T R_\Delta D$,
since both sides vanish on the dead pair and $D$ commutes with $R_\Delta$ on the others, so $W_K'=DW_K$ works with $\ell$
arbitrary on the dead pair. Under ReLU, $\diag(\ell)\mathrm{ReLU}(Wx)=\mathrm{ReLU}(\diag(\ell)Wx)$ holds for every $\ell\ge0$,
zero entries included.

(c) For a linear or affine box the class is the linear or affine maps, so a non-affine modified box function is outside it,
which is the definition of failing box-local mergeability.

(d) is \thmref{Proposition}{V.2}(d).
\end{proof}

\paragraph{The corollaries after \thmref{Theorem}{VI.3}}
\emph{Weight-space methods} merge by $\mu=\rho$ (\S\ref{proof:VI.1}); their nonlinearities act on the parameters. If an
embedding and an LM head share one tensor, a merge that changes one side must change the other or break the tie, which is why
an adapter on one side is not merged faithfully when the tie is not maintained. \emph{(IA)$^3$ on keys and values} is
$\diag(\ell)W_Kx=(\diag(\ell)W_K)x$ and likewise for $W_V$ (R1). \emph{On the FFN}, applied to the input of $W_{\rm down}$, it is
$W_{\rm down}\diag(\ell)h=(W_{\rm down}\diag(\ell))h$ (R1); neither step involves $\sigma$. \emph{Into $W_{\rm up}$}: under ReLU,
R3 gives $\diag(\ell)\mathrm{ReLU}(W_{\rm up}x)=\mathrm{ReLU}(\diag(\ell)W_{\rm up}x)$ for $\ell\ge0$, while an entry $\ell_i<0$
cannot fuse, since $\ell_i\mathrm{ReLU}(w_i\cdot x)\le0$ and is negative somewhere, whereas $\mathrm{ReLU}(w_i'\cdot x)\ge0$. Under
SwiGLU or GeGLU, R4 and R1 give \eqref{eq:merge-swiglu} for every $\ell$. \emph{Under GELU}, suppose
$\ell_i\,\mathrm{GELU}(w_i\cdot x)=\mathrm{GELU}(w_i'\cdot x)$ for all $x$, with $w_i\ne0$ so that the pre-activation ranges over
$\R$. If $\ell_i=0$, $w_i'=0$ works. If $\ell_i\ne0$, the right side vanishes on the hyperplane $w_i\cdot x=0$, and
$\mathrm{GELU}(t)=t\,\Phi_{\mathcal N}(t)$ vanishes only at $t=0$, so $w_i'\cdot x=0$ there and $w_i'=cw_i$. Then
$\ell_i\,\mathrm{GELU}(t)=\mathrm{GELU}(ct)$ for all $t$, and comparing the coefficients $a_1=\tfrac12$ and $a_2=1/\sqrt{2\pi}$
(\auxref{B.10}(c)) gives $c=\ell_i$ and $\ell_i=\ell_i^2$, so $\ell_i=1$. Hence a fused $W_{\rm up}'$ exists iff every
$\ell_i\in\{0,1\}$. \emph{A query scaling} fuses into $W_Q$ by R1; onto $W_K$ it needs R5's slide, hence $\ell$ constant on the
rotary pairs under RoPE (necessary under the genericity hypotheses of (b)), and with grouped-query attention one $W_K$ serves
several query heads $h$, so the same argument with full-rank $W_Q^{(h)}$ forces $W_K'=D_hW_K$ for every $h$ in the group, hence
$D_h$ constant across the group when $W_K$ has full row rank. \emph{SSF} after a box $Wx+b$ gives
$s\odot(Wx+b)+t=(\diag(s)W)x+(s\odot b+t)$, which needs a bias slot when $t\ne0$. \emph{LN-tuning} changes existing parameters, so
$\mu$ is the new parameter vector. \emph{Folding a LayerNorm} into the following boxes in a pre-LN block, where the norm feeds only
those boxes, is R7 followed by R6, which needs a bias slot for $W\beta$; an RMSNorm gain has no $\beta$. Folding a gain into a
tied embedding or head changes the other side of the tie. \emph{Not box-locally mergeable}: for $\sigma$ non-affine and nonzero
$u$, $d$, the map $y\mapsto y+u\,\sigma(d\T y)$ is non-affine, so a serial Houlsby or Pfeiffer adapter after a box $W$ with
$d\T W\ne0$, and a parallel adapter $x\mapsto Wx+u\,\sigma(d\T x)$, modify the box function non-affinely for some parameter
value; a routed mixture $x\mapsto Wx+\sum_eg_e(x)B_eA_ex$ with a softmax gate is real-analytic in its input and parameters and
non-affine at a witness (two experts with $B_1A_1=I$, $B_2A_2=0$ and gate logits differing by $\gamma\T x$, $\gamma\ne0$, give
$Wx+g(\gamma\T x)\,x$ with $g$ the logistic function, non-affine along any line on which $\gamma\T x$ varies), hence non-affine for
generic parameters (\auxref{B.10}), and linear, so mergeable, when the gate is constant; AdapterFusion with a constant gate is a fixed combination of nonlinear adapters,
still non-affine; LST adds the output of a nonlinear side network; and prefix-type methods are treated in
\thmref{Theorem}{VI.5}(f). In each case (c) applies. For LLaMA-Adapter with a two-token prompt $p_1,p_2$ and gate $g\ne0$,
the prompt pairs are $k_i=W_kp_i+b_k$ and $v_i=W_vp_i+b_v$, and a head on a length-one context computes
$W_o(W_vx+b_v)+g\,W_o\mathrm{Attn}_{q(x)}(P)+b_o$ with
$\mathrm{Attn}_{q}(P)=v_2+\operatorname{logit}^{-1}\big(\langle q,W_k\delta\rangle/\sqrt d\big)\,W_v\delta$, $\delta=p_1-p_2$.
If $W_q\T W_k\ne0$ and $W_oW_v\ne0$, choose $u$ and $\delta$ with $\langle W_qu,W_k\delta\rangle\ne0$ and $W_oW_v\delta\ne0$
(each condition fails only on a proper subspace or a nowhere-dense set). Along $x=tu$ the logistic factor is a bounded,
non-constant real-analytic function of $t$, hence not affine, so the head is not affine. By \auxref{B.10} the same holds for
generic prompts with two or more tokens, which is the entry for LLaMA-Adapter in the summary table of
Exposé~\ref{sec:merge}; with one prompt token, $\mathrm{Attn}_q(P)$ is constant and the head stays affine.

% ====================================================================================================
\subsection{Proposition VI.4}\label{proof:VI.4}

\begin{restate}{Proposition}{VI.4}{linear insertions, and ReFT}
(a) A linear adapter in parallel with a single linear box, $W_0x+sUDx$, \emph{is} LoRA: it has the same $\rho$. A linear
adapter parallel to a whole sublayer (He et al.'s parallel adapter around attention or the FFN) edits the weightless skip
connection rather than a base box, and is not box-locally mergeable when the sublayer is non-affine.

(b) Let $W_0\in\R^{m\times n}$ be bias-free (as LLaMA-type $W_o$ and $W_{\mathrm{down}}$ are), and $r\ge1$. A serial linear
adapter after $W_0$, $(I+UD)W_0$ with $U\in\R^{m\times r}$ and $D\in\R^{r\times m}$ (scale absorbed), has image
$W_0+\{Y\in\Mr:\operatorname{row}Y\subseteq\operatorname{row}W_0\}$. It is expressively equivalent to LoRA iff $\rk W_0=n$,
and weakly (unpointed) isomorphic to LoRA iff $W_0$ is invertible. The isomorphism $(U,D)\mapsto(U,DW_0)$ sends a base point
$(U_0,D_0)$ to $(U_0,D_0W_0)$. For a fixed $\mathrm T_A$ pointing $(0,A_0)$ with $\rk A_0=r$ and $1\le r<m=n$, a pointed
isomorphism exists iff $W_0$ is invertible and $\operatorname{row}A_0=\operatorname{row}(D_0W_0)$, and then $U_0=0$
automatically. With a bias $b_0$, the serial adapter also moves the bias by $UDb_0$, so on the affine box its image differs
from weight-only LoRA's. Houlsby's near-identity initialisation (weights from a zero-mean Gaussian with standard
deviation $10^{-2}$, truncated at two standard deviations) makes practical serial adapters slightly unpointed; the pointed
statement applies to zero-up initialisations.

(c) LoReFT is $\Phi(h)=h+R^\top(Wh+b-Rh)$, with $R\in\R^{r\times d}$ having orthonormal rows, applied at selected positions.
On one hidden vector it is the extension by an \textbf{identity box} followed by a rank-$\le r$ affine edit,
\[\Phi(h)=\big(I+R^\top(W-R)\big)h+R^\top b,\]
which is LoRA$_r$ on the \emph{affine} identity box $[I_d\mid0]$ in homogeneous coordinates, with $B=R^\top$ constrained to the
Stiefel manifold and $A=[W-R\mid b]$. Its image is $(I,0)+\MM_{\le r}(d\times(d+1))$, of dimension $r(2d+1-r)$. The bias is
confined to $\operatorname{col}B$, so for $r<d$ this is not LoRA$_r{}\oplus{}$BitFit, whose dimension is larger by $d-r$. It is
simulated, through $h(R,W,b)=(R^\top,[W-R\mid b])$, by LoRA$_r$ on that affine box pointed at $(R_0^\top,0)$, a $\mathrm T_B$
pointing with $\operatorname{col}B_0=\operatorname{row}R_0$; this homogeneous-coordinate LoRA has a bias column in $A$, unlike
HF LoRA. It is not simulated by $\mathrm T_A$-pointed LoRA$_r$, since its $\dim S_1=r(d+1)$ exceeds $rd$. It is not isomorphic
to LoRA$_r$: its parameter count is smaller by $r(r+1)/2$, and its residual symmetry is $\Ort(r)$, acting by
$(R,W,b)\mapsto(OR,OW,Ob)$. It is pointed iff $W=R$ and $b=0$. The reference implementation (pyreft) initialises $W$ and $b$
by the default \texttt{nn.Linear} init, independently of $R$, so practical LoReFT has a pointing defect
$R^\top((W_0-R_0)h+b_0)$: at initialisation it replaces the $R$-component of $h$ by a random affine function, and the
simulation above is then only a weak morphism. pyreft also applies dropout to the whole output, $h$ included, so the
training-time map is not $\Phi$.
\end{restate}

\begin{proof}
(a) The adapted box computes $x\mapsto(W_0+sUD)x$, so $\rho(U,D)=W_0+sUD$ on $\R^{m\times r}\times\R^{r\times n}$, which is
LoRA's $\rho(B,A)=W_0+sBA$ under $(U,D)=(B,A)$, with the zero-up pointing $U=0$ \citep[Table~1]{he2021unified}. A linear adapter
around a whole sublayer $F$ gives $x\mapsto x+F(x)+sUDx$; the term $sUDx$ is added to the weightless skip connection, which no
box of the base network computes, so there is no base box whose function the edit modifies. For a non-affine sublayer such as
$F(x)=W_2\sigma(W_1x)$ with $\sigma$ non-affine, no base box both reads $x$ and writes linearly into the block's output: $W_1$'s
output passes through $\sigma$, and $W_2$ reads $\sigma(W_1x)$ rather than $x$. So no box of $F$ can absorb $sUDx$ within its
class. If $F(x)=Mx+c$ were affine, $x+(M+sUD)x+c$ would fuse into $M$.

(b) \emph{Image.} $\rho(U,D)=(I+UD)W_0=W_0+UDW_0$. Each $UDW_0$ has rank at most $r$ and its rows are combinations of the rows
of $W_0$. Conversely, if $\rk Y\le r$ and $\operatorname{row}Y\subseteq\operatorname{row}W_0$, then $Y=YW_0^+W_0$, since
$W_0^+W_0$ is the orthogonal projector onto $\operatorname{row}W_0$; $X=YW_0^+$ has rank at most $r$, so $X=UD$ with
$U\in\R^{m\times r}$, $D\in\R^{r\times m}$, and $Y=UDW_0$. \emph{Expressive equivalence.} LoRA's image is $W_0+\Mr$. If
$\operatorname{row}W_0=\R^n$ the two images agree. If not, pick $v\notin\operatorname{row}W_0$; then $e_1v\T\in\MM_{\le1}\subseteq
\Mr$ (here $r\ge1$) has a row outside $\operatorname{row}W_0$, so it lies in LoRA's image minus $W_0$ but not in the
adapter's.

\emph{Weak isomorphism.} If $W_0$ is invertible, then $m=n$ and $h(U,D)=(U,DW_0)$ is a linear bijection
$\R^{m\times r}\times\R^{r\times m}\to\R^{m\times r}\times\R^{r\times n}$ with $W_0+U(DW_0)=\rho(U,D)$ (LoRA with $s=1$, or $U$
rescaled), so it is an isomorphism over $\Theta$; it sends $(U_0,D_0)$ to $(U_0,D_0W_0)$. Conversely, a weak isomorphism is a
diffeomorphism $h$ with $\rho_{\mathrm{LoRA}}\circ h=\rho$, so the images agree, which forces $\rk W_0=n$
(\thmref{Observation}{I.4}(c)), and the dimensions agree, $2mr=r(m+n)$, which forces $m=n$; so $W_0$ is square of full rank.

\emph{Pointed isomorphism.} Let $1\le r<m=n$, $\rk A_0=r$, and let the adapter be pointed at $(U_0,D_0)$, that is,
$U_0D_0W_0=0$. ``If'': $W_0$ invertible gives $U_0D_0=0$. Since $\operatorname{row}(D_0W_0)=\operatorname{row}A_0$ has dimension
$r$, $D_0W_0$ and hence $D_0$ have rank $r$, so $D_0$ has a right inverse and $U_0=0$. Two $r\times n$ matrices of rank $r$ with
the same row space differ by an invertible left factor, so $A_0=gD_0W_0$ for some $g\in\GL_r$. The map
$(U,D)\mapsto(Ug^{-1},gDW_0)$ is a linear bijection over $\Theta$ that sends $(U_0,D_0)=(0,D_0)$ to $(0,A_0)$. ``Only if'': a
pointed isomorphism is a weak one, so $W_0$ is invertible, and it preserves $S_1$ (\thmref{Observation}{I.6}). Differentiating,
\begin{align*}
S_1(\text{adapter at }(U_0,D_0))&=\{(XD_0+U_0Y)W_0\}=\{XD_0W_0+U_0Y'\},\\
S_1(\mathrm{LoRA}\text{ at }(0,A_0))&=\{XA_0\},
\end{align*}
with $X\in\R^{m\times r}$ and $Y'=YW_0$ ranging over $\R^{r\times n}$ because $W_0$ is invertible; and $\{XA_0\}=\{Z:
\operatorname{row}Z\subseteq\operatorname{row}A_0\}$. If $(U_0)_{ij}\ne0$, choose $v\notin\operatorname{row}A_0$, possible
since $r<n$, and $Y'=e_jv\T$: row $i$ of $U_0Y'$ is $(U_0)_{ij}v\T$, outside $\operatorname{row}A_0$, a contradiction. So $U_0=0$,
and $\{XD_0W_0\}=\{XA_0\}$ gives $\operatorname{row}(D_0W_0)=\operatorname{row}A_0$.

\emph{Bias.} $(I+UD)(W_0x+b_0)=(W_0+UDW_0)x+(b_0+UDb_0)$. On the affine box the adapter's image contains maps with bias
$b_0+UDb_0\ne b_0$ whenever $b_0\ne0$ (choose $UD$ with $UDb_0\ne0$), while weight-only LoRA keeps the bias $b_0$. \emph{Houlsby
initialisation.} With both factors drawn at random, $U_0D_0W_0\ne0$ almost surely, so the initial weight misses $W_0$.

(c) \emph{Algebra.} Expanding, $\Phi(h)=h+R\T(W-R)h+R\T b=\big(I+R\T(W-R)\big)h+R\T b$. In homogeneous coordinates $(h,1)$ an
affine map $h\mapsto Mh+c$ is the matrix $[M\mid c]\in\R^{d\times(d+1)}$, the identity box is $[I_d\mid0]$, and
\[\big[\,I+R\T(W-R)\ \big|\ R\T b\,\big]=[I_d\mid0]+R\T\,[W-R\mid b]=[I_d\mid0]+BA,\qquad B=R\T,\ A=[W-R\mid b].\]
\emph{Image.} As $(W,b)$ ranges freely, $V=[W-R\mid b]$ ranges over all of $\R^{r\times(d+1)}$, so the image is
$[I_d\mid0]+\{R\T V\}$ over Stiefel $R\T$. Every matrix of rank at most $r$ factors as $R\T V$ with $R\T$ having orthonormal
columns (\auxref{B.11}(a)), so the image is $[I_d\mid0]+\MM_{\le r}(d\times(d+1))$, of dimension $r(d+(d+1)-r)=r(2d+1-r)$
(\bgref{C.12}). The bias $R\T b$ lies in $\operatorname{col}R\T=\operatorname{col}B$. LoRA$_r{}\oplus{}$BitFit has image
$\{[I_d+Y\mid c]:\rk Y\le r,\ c\in\R^d\}$, of dimension $r(2d-r)+d$, which exceeds $r(2d+1-r)$ by $d-r$; for $r<d$ the images
differ. \emph{Simulation.} The map $h(R,W,b)=(R\T,[W-R\mid b])$ is smooth, and LoRA$_r$ on the affine box evaluated at $h(R,W,b)$ is
$[I_d\mid0]+R\T[W-R\mid b]=\Phi$, so $h$ is a morphism to that LoRA$_r$. It sends the neutral point $(R_0,R_0,0)$ to $(R_0\T,0)$, a pointing in the stratum $\mathrm T_B$
($A_0=0$, $B_0$ of rank $r$) with $\operatorname{col}B_0=\operatorname{row}R_0$. \emph{$S_1$.} At the neutral point
$[W-R\mid b]=0$, so the derivative of $R\T[W-R\mid b]$ is $R_0\T[\dot W-\dot R\mid\dot b]$, and as $(\dot W,\dot b)$ range freely
$S_1=\{R_0\T V\}$, of dimension $r(d+1)$ because $R_0\T$ is injective. LoRA$_r$ pointed at $(0,A_0)$ with $\rk A_0=r$ has
$S_1=\{XA_0\}$, of dimension $dr$. A pointed morphism would give an inclusion of $S_1$'s (\thmref{Observation}{I.6}), impossible
since $r(d+1)>rd$.

\emph{Not isomorphic.} The parameter space is the Stiefel manifold of orthonormal $r$-frames in $\R^d$, of dimension
$rd-r(r+1)/2$, times $\R^{r\times d}\times\R^r$, of total dimension $2rd+r-r(r+1)/2$; LoRA$_r$ on the affine box has dimension
$r(2d+1)$. The difference $r(r+1)/2$ is positive, so no diffeomorphism, hence no isomorphism, exists. \emph{Symmetry.} For
$O\in\Ort(r)$, $OR$ has orthonormal rows and $(OR)\T(OW-OR)=R\T(W-R)$, $(OR)\T Ob=R\T b$, so $\Phi$ is unchanged. Conversely, if
$R\T V=R'^{\top}V'$ with $\rk V=r$, then the column space of the product is $\operatorname{col}R\T=\operatorname{col}R'^{\top}$, so
$R'=OR$ for some $O\in\Ort(r)$, and injectivity of $R\T$ gives $V'=OV$, that is, $W'=OW$ and $b'=Ob$: the generic fibres are the
$\Ort(r)$-orbits, of dimension $r(r-1)/2$, against $r^2$ for LoRA. \emph{Pointedness.} $\Phi=\id$ iff $R\T[W-R\mid b]=0$ iff
$W=R$ and $b=0$, because $R\T$ is injective.

\emph{The pyreft initialisation.} At $(R_0,W_0,b_0)$ with $W_0$, $b_0$ drawn independently of $R_0$, $\Phi_0(h)-h=R_0\T\big((W_0-R_0)
h+b_0\big)$, which is nonzero almost surely. Since $R_0R_0\T=I_r$, $R_0\Phi_0(h)=R_0h+(W_0h+b_0-R_0h)=W_0h+b_0$, while
$(I-R_0\T R_0)\Phi_0(h)=(I-R_0\T R_0)h$: the $R_0$-component of $h$ is replaced by the random affine function $W_0h+b_0$ and the
rest is kept. The map $h$ sends this initial point to $(R_0\T,[W_0-R_0\mid b_0])$, not to the base point $(R_0\T,0)$, so it is a
weak morphism only. The remark on dropout describes the implementation \citep{wu2024reft}.
\end{proof}

% ====================================================================================================
\subsection{Theorem VI.5}\label{proof:VI.5}

\begin{restate}{Theorem}{VI.5}{attention is a homomorphism followed by a perspective map; known in part}
Fix a query $q$. Multisets of key--value pairs form a commutative monoid under $\uplus$. Define
\[\Sigma_q(K)=\Big(\sum e^{\langle q,k\rangle/\sqrt d},\ \sum e^{\langle q,k\rangle/\sqrt d}\,v\Big).\]

(a) $\Sigma_q$ is a monoid homomorphism into $(\R_{\ge0}\times\R^{d_v},+)$, landing in $\{Z>0\}\cup\{(0,0)\}$. On nonempty
multisets, $\mathrm{Attn}_q=\pi\circ\Sigma_q$ with the partial map $\pi(Z,S)=S/Z$, defined for $Z>0$.

(b) \textbf{Prefix is a gated parallel adapter.} For nonempty $P$ and $C$,
\[\mathrm{Attn}_q(P\uplus C)=(1-\lambda_q)\,\mathrm{Attn}_q(C)+\lambda_q\,\mathrm{Attn}_q(P),\qquad
\lambda_q=\frac{Z_q(P)}{Z_q(P)+Z_q(C)}.\]

(c) \textbf{Per-layer pattern rigidity.} Within a single layer, for fixed $q$ and a fixed content multiset $C$: the content
attention weights are multiplied by $1-\lambda_q$, their ratios are unchanged, and the output lies on the segment from
$\mathrm{Attn}_q(C)$ to $\mathrm{Attn}_q(P)\in\operatorname{conv}V_P$, where $V_P$ is the set of prefix values. This rigidity
is relative to the same layer's computation without the prefix; it is relative to the pretrained pattern only at the lowest
prefixed layer. Lower-layer prefixes change the queries and the content keys and values of higher layers, so attention
patterns at the level of the network are \emph{not} rigid.

(d) \textbf{LLaMA-Adapter} normalises $P$ and $C$ separately and adds a zero-initialised gate. It computes
$\mathrm{Attn}_q(C)+\tanh(g_l)\,\mathrm{Attn}_q(P)$, with one gate per head in the paper; HF PEFT uses one raw
zero-initialised scalar per layer, without $\tanh$. It uses two perspective maps, leaves the content weights exactly
unchanged, and is pointed at $g=0$. Its prompt keys and values are the frozen $W_k,W_v$ applied to one learned prompt per
adapted layer, and it acts in the top layers only.

(e) \textbf{Prompt $\to$ deep prefix in causal decoders.} Assume that the soft prompt comes first, that attention is causal,
and that content positions are offset by the prompt length $m$ in both methods (as HF PEFT does via \texttt{cache\_position}).
Then at every layer the states at prompt positions depend only on $p$, and $h:p\mapsto(K_j(p),V_j(p))_j$ satisfies
\[\Phi\circ\rho_{\mathrm{prefix}}\circ h=\Phi\circ\rho_{\mathrm{prompt}}\]
as functions of the content tokens, read at content positions. The two methods live on different extensions, so this is an
arrow of the unpointed slice over the realisation $C^\infty(X_{\mathrm{content}},Y_{\mathrm{content}})$, not of $\Meth$.
Under RoPE without the offset, a different $h$ works: rotate each prefix key by $-m$ (positions $-m,\dots,-1$). For additive
relative biases (ALiBi, T5) the answer depends on how the implementation places the prefix keys. The construction does not
apply in bidirectional encoders of depth $\ge2$.

(f) \textbf{Obstructions.} Call a head \emph{standard} if it is vanilla softmax attention without torch-style
\texttt{add\_bias\_kv} (which would absorb a one-token prefix), \texttt{add\_zero\_attn}, learned attention sinks or
softmax-plus-one. On length-one contexts, every weight setting of a standard head is affine in $x$. Some prefix makes the
prefixed head non-affine on length-one contexts iff $W_o\ne0$ and ($W_q\ne0$ or $W_k^\top b_q\ne0$); for bias-free heads, iff
$W_o\ne0$ and $W_q\ne0$. For such frozen weights and generic prefixes the prefixed head is not affine, because the gate
$\lambda(x)$ is not constant. So prefix tuning and P-tuning v2 are not box-locally mergeable into a standard head. When
$W_oW_v\ne0$, no prefix is neutral, so prompt and prefix tuning are unpointed.
\end{restate}

\begin{proof}
Write $\epsilon_q(k)=e^{\langle q,k\rangle/\sqrt d}>0$, and for a multiset $K$ put $Z_q(K)=\sum_{(k,v)\in K}\epsilon_q(k)$ and
$S_q(K)=\sum_{(k,v)\in K}\epsilon_q(k)v$, so $\Sigma_q(K)=(Z_q(K),S_q(K))$.

(a) Sums over a multiset union split, $\Sigma_q(K\uplus K')=\Sigma_q(K)+\Sigma_q(K')$, and $\Sigma_q(\emptyset)=(0,0)$, so
$\Sigma_q$ is a monoid homomorphism. Each term of $Z_q$ is positive, so $Z_q(K)>0$ for nonempty $K$ and $Z_q(\emptyset)=0$. By
definition, $\mathrm{Attn}_q(K)=S_q(K)/Z_q(K)=\pi(\Sigma_q(K))$ for nonempty $K$; this $\pi$ is the perspective map of
\bgref{D.34}.

(b) With $\lambda=Z_P/(Z_P+Z_C)$ and $Z_P,Z_C>0$,
\[\pi(Z_P+Z_C,S_P+S_C)=\frac{S_P+S_C}{Z_P+Z_C}=\frac{Z_P}{Z_P+Z_C}\frac{S_P}{Z_P}+\frac{Z_C}{Z_P+Z_C}\frac{S_C}{Z_C}
=\lambda\,\mathrm{Attn}_q(P)+(1-\lambda)\,\mathrm{Attn}_q(C).\]

(c) In $P\uplus C$ a content pair $(k,v)\in C$ has weight $\epsilon_q(k)/(Z_P+Z_C)=(1-\lambda)\,\epsilon_q(k)/Z_C$, which is
$1-\lambda$ times its weight in $C$ alone; so all content weights scale by one factor and their ratios are unchanged. By (b)
the output lies on the segment from $\mathrm{Attn}_q(C)$ to $\mathrm{Attn}_q(P)$, with $\lambda\in(0,1)$, and
$\mathrm{Attn}_q(P)=\sum_{(k,v)\in P}(\epsilon_q(k)/Z_P)\,v$ is a convex combination of the prefix values. At the lowest
prefixed layer the query and the content pairs are computed from inputs that no prefix has touched, so they are those of the
pretrained model, and the content weights are the pretrained ones times $1-\lambda_q$. At a higher layer the inputs come from
lower layers whose outputs (b) has changed, so $q$ and $C$ themselves differ from the pretrained ones, and the statement
compares only with the same layer's computation on those inputs. The failure of network-level rigidity is shown in
\citet[\S6]{petrov2023prompting}, and universality of prefix-tuning in \citet{petrov2024prompting}.

(d) LLaMA-Adapter \citep{zhang2023llama} computes the two softmaxes separately, so its output is
$\pi(\Sigma_q(C))+\tanh(g_l)\,\pi(\Sigma_q(P))$, with two perspective maps. A content weight is $\epsilon_q(k)/Z_q(C)$, the same as
without the adapter. At $g=0$, $\tanh(0)=0$ (and HF's raw gate is $0$), so the output is $\mathrm{Attn}_q(C)$, the pretrained
one, and the method is pointed there. The remaining clauses describe how the method is built.

(e) Number the positions of the prompt-tuned model $0,\dots,m-1$ for the soft prompt $p=(p_0,\dots,p_{m-1})$ and $m,\dots,m+n-1$
for the content, and give the content the same positions $m,\dots,m+n-1$ in the prefix-tuned model. Let $K_j(p),V_j(p)$ be the
keys and values, after the positional encoding, that the prompt-tuned model forms at the prompt positions in layer $j$, as they
are cached. We show by induction over layers $j$ that (i) the hidden states of the prompt-tuned model at prompt positions
depend only on $p$, and (ii) the hidden states at content positions agree in the prompt-tuned model and in the prefix-tuned
model with prefix $h(p)=(K_j(p),V_j(p))_j$. At the input, (i) holds trivially and (ii) holds because both models embed the same
content tokens at the same positions. Suppose (i) and (ii) hold at the input of layer $j$. Under causal masking a prompt
position attends only to prompt positions, so its output depends only on prompt states, which gives (i) at the next layer. A
content position $t$ attends to the prompt positions and to the content positions $u\le t$. Its query, and the content keys
and values, are computed from equal content states with equal positional encodings, by (ii). Its prompt keys and values are
$K_j(p),V_j(p)$ in the prompt-tuned model by definition, and they are supplied verbatim by the prefix in the other. With
rotary encodings, a score is $\langle R_tq,R_sk_s\rangle$ in both models, the cached key $R_sk_s$ being the same object. So the
query and the multiset of key--value pairs are the same, and the attention outputs agree; the position-wise sublayers then give
(ii) at the next layer. Reading the content positions at the top gives $\Phi\circ\rho_{\rm prefix}\circ h=\Phi\circ
\rho_{\rm prompt}$. The prompt-tuned model and the prefix-tuned model are methods on different extensions of the base network
(input tokens in one, key--value slots in the other), so $h$ is a map over the common realisation as functions of the content,
an arrow of the slice over $C^\infty(X_{\rm content},Y_{\rm content})$.

Without the offset, the prefix-tuned model places the content at positions $0,\dots,n-1$. For RoPE put
$h'(p)=\big(R_{-m}K_j(p),V_j(p)\big)_j$, rotating the cached key from prompt position $s$ to $s-m$. A content query at position
$t-m$ then scores $\langle R_{t-m}q,R_{s-m}k_s\rangle=\langle q,R_{s-t}k_s\rangle=\langle R_tq,R_sk_s\rangle$, content--content
scores depend only on differences of positions and are unchanged, and values carry no rotation, so the induction goes through
with $h'$. For additive relative biases the score of a content query against a prefix key depends on the distance the
implementation assigns to that key, and the induction goes through exactly when these distances equal $t-s$. In a bidirectional
encoder of depth at least $2$, prompt positions attend to content positions, so already the second layer's prompt states depend on
the content, and step (i) fails.

(f) \emph{Affine without prefix.} On a context of one token $x$ the softmax has one term, equal to $1$, so a standard head
computes $W_o(W_vx+b_v)+b_o$, which is affine in $x$ for every weight setting. (With \texttt{add\_zero\_attn}, a sink or
softmax-plus-one there is a second term, and this fails.)

\emph{With a prefix.} Let the prefix be $P=\{(k_i,v_i)\}$ with $N_P\ge1$ fixed pairs, and write $q(x)=W_qx+b_q$,
$k(x)=W_kx+b_k$, $v(x)=W_vx+b_v$, the content score $\varsigma(x)=\langle q(x),k(x)\rangle/\sqrt d$ and
$Z_P(x)=\sum_i\epsilon_{q(x)}(k_i)$. By (b),
\begin{equation}\label{eq:C3-prefix-head}
\mathrm{head}_P(x)=W_o\big[(1-\lambda(x))\,v(x)+\lambda(x)\,\mathrm{Attn}_{q(x)}(P)\big]+b_o,\qquad
\lambda(x)=\frac{Z_P(x)}{Z_P(x)+e^{\varsigma(x)}}\in(0,1).
\end{equation}
\emph{Only if.} If $W_o=0$ the output is the constant $b_o$. If $W_q=0$ and $W_k\T b_q=0$, then $q=b_q$ is constant, every prefix
score $\langle b_q,k_i\rangle/\sqrt d$ is constant, and the content score $\langle b_q,W_kx+b_k\rangle/\sqrt d=\big((W_k\T b_q)\T
x+\langle b_q,b_k\rangle\big)/\sqrt d$ is constant; so $\lambda$ and $\mathrm{Attn}_q(P)$ are constant and the output is affine.

\emph{If.} Suppose $W_o\ne0$ and ($W_q\ne0$ or $W_k\T b_q\ne0$). Take all $N_P$ prefix pairs equal to $(k_1,v_1)$; then
$\mathrm{Attn}_{q(x)}(P)=v_1$ and $\lambda(x)=\operatorname{logit}^{-1}\big(\delta(x)\big)$ with
$\delta(x)=\langle q(x),k_1\rangle/\sqrt d+\log N_P-\varsigma(x)$. Restrict to a line $x=tu$. If $W_q\ne0$, choose $u$ with
$W_qu\ne0$; $\varsigma(tu)$ is a fixed polynomial of degree at most two in $t$, and $\langle q(tu),k_1\rangle$ has the slope
$\langle W_qu,k_1\rangle$, which $k_1$ can make any real number, so $k_1$ can be chosen with $\delta(tu)$ non-constant. If $W_q=0$ and
$W_k\T b_q\ne0$, choose $u$ with $\langle W_k\T b_q,u\rangle\ne0$; then $\varsigma(tu)$ has nonzero slope while the prefix score is
constant, so $\delta(tu)$ is non-constant for every $k_1$. In both cases $\lambda(t):=\lambda(tu)$ is real-analytic, takes values
in $(0,1)$ and is non-constant, so $\lambda''(t_0)\ne0$ for some $t_0$; otherwise $\lambda$ would be affine and bounded, hence
constant. Choose $v_*$ with $W_ov_*\ne0$ and put $v_1=Mv_*$. Since $v(tu)$ is affine in $t$,
\[\frac{d^2}{dt^2}\mathrm{head}_P(tu)=\lambda''(t)\big(MW_ov_*-W_ov(tu)\big)-2\lambda'(t)\,W_oW_vu ,\]
and at $t_0$ the term $M\lambda''(t_0)W_ov_*$ dominates for large $M$, so the second derivative is nonzero and the head is not
affine. For bias-free heads $W_k\T b_q=0$, which gives the second form of the criterion.

\emph{Generic prefixes.} The prefixed output is real-analytic in $(x,P)$ on the connected space of inputs and prefixes with $N_P$
pairs. It is affine in $x$ for a given $P$ iff all second partial derivatives in $x$ vanish at every $x$; this is the common zero
set, over $x$ and pairs of coordinates, of real-analytic functions of $P$. The witness above is a prefix at which one of these
functions is nonzero, so that function is not identically zero, and the prefixes with affine output lie in its zero set, which
is closed, nowhere dense and null (\auxref{B.10}). The same argument applied to $\lambda$ shows that the gate is non-constant
for generic prefixes. Box-locally, the class of a standard head is the set of functions it computes with some weights; on
length-one contexts these are all affine, while for generic prefixes the prefixed head is not, so prefix tuning and P-tuning v2
are not box-locally mergeable into the head.

\emph{No neutral prefix.} A neutral prefix would leave the head unchanged on every input, in particular on length-one contexts,
where by \eqref{eq:C3-prefix-head} this means $\lambda(x)\big(W_o\mathrm{Attn}_{q(x)}(P)-W_ov(x)\big)=0$, and since $\lambda(x)>0$,
$W_o\mathrm{Attn}_{q(x)}(P)=W_o(W_vx+b_v)$ for all $x$. The left side lies in the compact set $W_o\operatorname{conv}V_P$; if
$W_oW_v\ne0$, choose $x_0$ with $W_oW_vx_0\ne0$, and along $x=tx_0$ the right side is unbounded. So no prefix is neutral. For a
base network consisting of one such head, the extension by prefixes therefore has no neutral point, and by (e), at the first layer
of a causal decoder a soft prompt acts through its key--value pairs exactly as a one-layer prefix, so the same base network admits
no neutral prompt: prompt and prefix tuning are unpointed.
\end{proof}

\begin{remark*}{numerical checks}
The values reported after \thmref{Theorem}{VI.5} (prompt-to-prefix agreement $9\times10^{-16}$ with RoPE and the offset, $1.3$
without it and $2\times10^{-15}$ with keys rotated by $-m$; an affine head with $W_o=0$ although $W_q\T W_k\ne0$, and a non-affine
one with $W_q\ne0$ and $W_q\T W_k=0$) instantiate parts (e) and (f) as proved above: the criterion of (f) involves $W_o$, $W_q$ and
$W_k\T b_q$, not $W_q\T W_k$ (Appendix~\ref{app:repro}).
\end{remark*}
